\documentclass[11pt]{amsart}
\usepackage{amssymb,amsmath,amsthm}
\usepackage{mathtools}
\usepackage{booktabs}
\usepackage[margin=1.15in]{geometry}
\usepackage[colorlinks=true,linkcolor=blue,citecolor=blue,urlcolor=blue]{hyperref}

\newtheorem{theorem}{Theorem}[section]
\newtheorem{lemma}[theorem]{Lemma}
\newtheorem{proposition}[theorem]{Proposition}
\newtheorem{corollary}[theorem]{Corollary}
\newtheorem{conjecture}[theorem]{Conjecture}
\theoremstyle{definition}
\newtheorem{definition}[theorem]{Definition}
\theoremstyle{remark}
\newtheorem{remark}[theorem]{Remark}
\numberwithin{equation}{section}
\setlength{\emergencystretch}{3em}

\newcommand{\N}{\mathbb{N}}
\newcommand{\R}{\mathbb{R}}
\newcommand{\diff}{\,\mathrm{d}}
\newcommand{\usat}{u_{\mathrm{sat}}}
\newcommand{\Sp}{\mathcal{S}_p}
\newcommand{\cE}{\mathcal{E}}
\newcommand{\cJ}{\mathcal{J}}
\newcommand{\Proj}{\mathsf{P}}
\newcommand{\ia}{\mathfrak{i}}
\DeclareMathOperator{\osc}{osc}
\DeclareMathOperator{\arc}{arc}
\DeclareMathOperator{\sgn}{sgn}
\DeclareMathOperator{\dist}{dist}
\DeclareMathOperator{\Var}{Var}
% PLACEHOLDER (decision D2 of the series plan): the hard-edge tail bound is to be proved in Paper II; its number is
% not yet fixed. Every use in this paper goes through this macro. See NOTES.md, "Placeholders".
\newcommand{\tailbound}{\cite[Proposition~(tail bound)]{PaperII}}

\title[Stability and instability of recurrence asymptotics]{Stability and instability of recurrence asymptotics for
discrete measures with power-law atoms}
\author{Gordon [Surname]}
\date{Draft of 15 October 2026}

\subjclass[2020]{Primary 42C05; Secondary 47B36, 26A12, 11M06}
\keywords{Orthogonal polynomials, Jacobi matrices, recurrence coefficients, discrete measures, slowly varying
functions, Christoffel transformation, Hankel determinants, zeta values}

\begin{document}

\begin{abstract}
We continue the study of the recurrence coefficients $a_n,b_n$ of discrete measures $\sum_kw_k\delta_{k^{-p}}$ with
power-law atoms. For weights $w_k=k^{-s}\ell(k)$ with $s>1$, a previous paper proved the pointwise law
$n^pb_n\to c(p)$, $n^pa_n\to2c(p)$, with rate $O(1/n)$, when $\ell$ is analytic and slowly varying. Here we show by
soft methods that the law holds for every slowly varying $\ell$, bounded or not, with a rate controlled by how well
$\ell$ can be smoothed in the logarithmic variable. Slow variation cannot be replaced by boundedness: for every
$p>0$, $s>1$ and every small $\varepsilon>0$ there are weights $k^{-s}e^{\varepsilon_k}$ with
$|\varepsilon_k|\le\varepsilon$ for which $n^pb_n$ does not converge, and others for which $n^pa_n$ does not
converge, although the averaged law of the first paper holds for all of them. For pure powers we then let $s$ vary.
A Christoffel identity gives
$\sum_{m\le n}\log(b_m(s+p)/b_m(s))+\frac p2\log n\to\frac12\log(c(p)\zeta(s)/\zeta(s+p))$. A variational argument
extends it to
$\sum_{m\le n}\log(b_m(s_1)/b_m(s_0))=-\frac{s_1-s_0}{2}\log n+\frac12\log(c(p)^{(s_1-s_0)/p}\zeta(s_0)/\zeta(s_1))
+o(1)$ for all $s_0,s_1>1$, and gives the law uniformly on compact sets of exponents (for $p\le1$, inside an
explicit window). At $p=1$ this determines the
averaged $1/n$ coefficient $-(s-1)/2$ for every $s>1$. All proofs rest on the first-variation formula for $\log h_n$
under multiplicative changes of the weights, combined with localisation of the spectral weights.
\end{abstract}

\maketitle

% Section 1: Introduction.
\section{Introduction}\label{sec:intro}

\subsection{Setting and results}\label{sec:results}
For $p>0$, $s>1$ and a sequence $\ell\colon\N\to(0,\infty)$ consider the discrete measure
\begin{equation}\label{eq:sigma}
  \sigma_{p,s,\ell}:=\sum_{k\ge1}k^{-s}\ell(k)\,\delta_{k^{-p}},\qquad \sigma_{p,s}:=\sigma_{p,s,1}.
\end{equation}
Its atoms $x_k=k^{-p}$ accumulate only at $0$. Let $P_n$ be its monic orthogonal polynomials and
$h_n:=\int P_n^2\diff\sigma$. They satisfy the three-term recurrence
\begin{equation}\label{eq:rec}
  xP_n=P_{n+1}+a_nP_n+b_n^2P_{n-1}\qquad(n\ge0,\ b_0:=0,\ b_n>0),
\end{equation}
so that $b_n^2=h_n/h_{n-1}$, and $a_n,b_n$ are the entries of the Jacobi matrix of the measure. Put
\begin{equation}\label{eq:cp}
  I(p):=\frac{\sqrt\pi}{2}\,\frac{\Gamma(\frac1p+\frac12)}{\Gamma(\frac1p+1)},\qquad
  c(p):=\frac14\Bigl(\frac{\pi}{2I(p)}\Bigr)^{p},
\end{equation}
so that $c(1)=\frac12$ and $c(2)=\pi^2/16$. We say that a measure satisfies the \emph{pointwise law} if
\begin{equation}\label{eq:law}
  n^pb_n\longrightarrow c(p)\qquad\text{and}\qquad n^pa_n\longrightarrow2c(p)\qquad(n\to\infty).
\end{equation}
The first paper of this series \cite{PaperI} proves \eqref{eq:law} in a doubly averaged form for all weights
between $c_1k^{-s}$ and a summable sequence; the constant $c(p)$ comes from an explicitly solvable constrained
equilibrium problem. The second paper \cite{PaperII} proves \eqref{eq:law} itself, with rate $O(1/n)$, by
Riemann--Hilbert analysis, when $\ell$ extends to a slowly varying analytic function in a sector, and in particular
for the pure powers $\sigma_{p,s}$. The present paper uses the second paper only through the numbered statements
restated in \S\ref{sec:imports}, and contains no Riemann--Hilbert analysis. It asks which weights the pointwise law
tolerates, and how the coefficients of the pure powers depend on $s$.

\begin{definition}[slow variation]\label{def:sv}
A sequence $\ell\colon\N\to(0,\infty)$ is \emph{slowly varying} if, for every $\lambda\ge1$,
\[
  \osc_n(\lambda):=\max\bigl\{|\log\ell(k)-\log\ell(n)|:\ k\in\N,\ n/\lambda\le k\le\lambda n\bigr\}\longrightarrow0
  \qquad(n\to\infty).
\]
\end{definition}

The case $\lambda=2$ implies the others (Remark~\ref{rem:lambda2}). No boundedness, monotonicity or measurability is
involved.

\begin{theorem}[every slowly varying factor]\label{thm:sv}
Let $p>0$, $s>1$, and let $\ell$ be slowly varying. Then the recurrence coefficients of $\sigma_{p,s,\ell}$ satisfy
$n^pb_n\to c(p)$ and $n^pa_n\to2c(p)$. More precisely,
\begin{equation}\label{eq:sv-rate}
  |n^pb_n-c(p)|+|n^pa_n-2c(p)|=O\bigl(1/n+\rho_\ell(\lceil n/M_0\rceil)\bigr),
\end{equation}
where $\rho_\ell(m)\to0$ is the error of Gaussian smoothing of $\log\ell$ in the logarithmic variable
(Definition~\ref{def:smoothing}) and $M_0=M_0(p)\ge4$ is an integer depending only on $p$ (Lemma~\ref{lem:saturation}).
\end{theorem}

For example, $\ell(k)=\exp\sqrt{\log k}$ and $\ell(k)=\exp(\sqrt{\log k}\,\sin\log\log k)$ ($k\ge3$) are covered; the
second is not eventually comparable with any power of $\log k$. The only analytic input is the case of analytic
factors, \cite[Theorem~1.2]{PaperII}, which we apply to a smoothed factor.

Slow variation cannot be replaced by boundedness. To state this, put $\usat(p):=(4c(p))^{-1/p}$; we show
$\usat(p)<1$ in Lemma~\ref{lem:usat}.

\begin{theorem}[bounded rough weights]\label{thm:rough}
Let $p>0$ and $s>1$.
\begin{enumerate}
\item[\textup{(i)}] Let $0<\tau_1<\tau_2$ with $\tau_2>\usat(p)$; every $\tau_2\ge1$ qualifies, for example
  $[\tau_1,\tau_2]=[0.6,1.6]$. There is $\gamma_W>0$, depending only on $p,\tau_1,\tau_2$ and given by an integral
  over the arcsine law \eqref{eq:window-constants}, with the following property. For every
  $\varepsilon\in(0,\gamma_W/16]$ there are integers $N_1<N_2<\cdots$ and numbers
  $\varepsilon_k\in\{-\varepsilon,0,\varepsilon\}$, vanishing outside $\bigcup_i[\tau_1N_i,\tau_2N_i]$, such that the
  weights $w_k=k^{-s}e^{\varepsilon_k}$ satisfy
  \[
    \limsup_{n\to\infty}\log\frac{b_n^2}{b_{n-1}^2}\ge\frac{\varepsilon\gamma_W}{4}.
  \]
\item[\textup{(ii)}] Let $\ia\in(0,1]$ be the number defined in \eqref{eq:ia} and $\gamma_{\mathrm a}:=c(p)\ia/2$.
  For every $\varepsilon\in(0,\ia/16]$ there are $\varepsilon_k\in\{-\varepsilon,0,\varepsilon\}$ such that the weights
  $w_k=k^{-s}e^{\varepsilon_k}$ satisfy
  \[
    \limsup_{n\to\infty}\bigl(n^pa_n-2c(p)\bigr)\ge\frac{\varepsilon\gamma_{\mathrm a}}2,\qquad
    \liminf_{n\to\infty}\bigl(n^pa_n-2c(p)\bigr)\le-\frac{\varepsilon\gamma_{\mathrm a}}2.
  \]
\item[\textup{(iii)}] For all weights with $e^{-\varepsilon}k^{-s}\le w_k\le e^{\varepsilon}k^{-s}$, in particular those
  of \textup{(i)} and \textup{(ii)}, the Hankel determinants and the coefficients $b_n$ of $\sum_kw_k\delta_{k^{-p}}$
  satisfy the averaged law of \cite{PaperI}, \eqref{eq:paperI-logD} and \eqref{eq:paperI-cesaro} below. In case
  \textup{(i)} the sequence $n^pb_n$ has no limit in $[0,\infty]$; in case \textup{(ii)} the sequence $n^pa_n$ has no
  limit.
\end{enumerate}
\end{theorem}

Numerically $\ia\approx0.314$, while the proved constant $\gamma_W$ is very small (about $10^{-7}$ for $p=2$ and
$[0.6,1.6]$); the effect itself is of size about $\varepsilon$ (Remark~\ref{rem:constants} and
Section~\ref{sec:numerics}). The only analytic input of Theorem~\ref{thm:rough} is the pure-power law without rate,
\cite[Corollary~1.3]{PaperII}.

For pure powers we write $a_n(s),b_n(s),h_n(s)$ for the quantities of $\sigma_{p,s}$, and $\zeta$ for the Riemann zeta
function, so that $\sigma_{p,s}$ has total mass $\zeta(s)$. Some of the statements below use a bound for the spectral
weights of $\sigma_{p,s}$ near the accumulation point of the atoms, which follows from \tailbound{} and an elementary
lemma (Lemma~\ref{lem:allM}) for $s$ in the set
\begin{equation}\label{eq:Sp}
  \Sp:=(1,1+2p)\quad\text{if }0<p\le1,\qquad \Sp:=(1,\infty)\quad\text{if }p>1.
\end{equation}

\begin{theorem}[pure powers as $s$ varies]\label{thm:s}
Let $p>0$.
\begin{enumerate}
\item[\textup{(i)}] For every $s>1$,
  \[
    \sum_{m=1}^n\log\frac{b_m(s+p)}{b_m(s)}+\frac p2\log n\longrightarrow\frac12\log\frac{c(p)\zeta(s)}{\zeta(s+p)}.
  \]
\item[\textup{(ii)}] For all $s_0,s_1>1$,
  \[
    \sum_{m=1}^n\log\frac{b_m(s_1)}{b_m(s_0)}=-\frac{s_1-s_0}2\log n
    +\frac12\log\frac{c(p)^{(s_1-s_0)/p}\zeta(s_0)}{\zeta(s_1)}+o(1).
  \]
  Equivalently $h_n(s_1)/h_n(s_0)=c(p)^{(s_1-s_0)/p}n^{-(s_1-s_0)}(1+o(1))$, and this holds uniformly for $s_0,s_1$
  in compact subsets of $\Sp$.
\item[\textup{(iii)}] For every compact $K\subset\Sp$, $\sup_{s\in K}|n^pb_n(s)-c(p)|\to0$. If moreover $p>1$, or
  $0<p\le1$ and $1+p\notin K$, then also $\sup_{s\in K}|n^pa_n(s)-2c(p)|\to0$.
\end{enumerate}
\end{theorem}

Part (i) uses only the limits of \cite[Corollary~1.3]{PaperII}; parts (ii) and (iii) also use \tailbound{}. Summing
the $1/n$ expansion that a pointwise statement would give, part (ii) says that, in logarithmic Ces\`aro mean, the
coefficient of $1/n$ in $n^pb_n(s)/c(p)$ decreases in $s$ with slope exactly $-\frac12$. An anchor fixes it at $p=1$.

\begin{corollary}[the averaged $1/n$ coefficient at $p=1$]\label{cor:p1}
For every $s>1$ there is $C(s)\in\R$ such that the coefficients of $\sigma_{1,s}=\sum_kk^{-s}\delta_{1/k}$ satisfy
\[
  \sum_{j=1}^n\log\bigl(2j\,b_j(s)\bigr)=-\frac{s-1}2\log n+C(s)+o(1).
\]
Equivalently, with $D_n(s):=\det\bigl(\zeta(s+i+j)\bigr)_{0\le i,j\le n-1}$,
\[
  \log\frac{D_{n+1}(s)}{D_n(s)}=\log\zeta(s)-2\log n!-2n\log2-(s-1)\log n+2C(s)+o(1).
\]
\end{corollary}

For $p=1$ and $s\in\frac32+\N_0$ the first statement is proved without the bound of \tailbound{}, and so is its
analogue for $p=2$ and $s\in2+2\N_0$ (Corollary~\ref{cor:lattice}). For $p>1$ the pointwise $1/n$ term is known,
$n^pb_n(s)=c(p)(1-\frac{s-1}{2n}+O(n^{-2}))$ \cite[Theorem~1.4]{PaperII}; its sum gives the logarithmic mean up to
$O(1)$, and Theorem~\ref{thm:s}(ii) adds the exact differences of the constants. For $0<p<1$ only differences in $s$
are determined (Section~\ref{sec:open}).

\subsection{Context}\label{sec:context}
\emph{Averaged and pointwise laws.} Specialised to the atoms $k^{-p}$, \cite[Theorem~1.1 and Corollary~1.2]{PaperI}
state that if $\sum_kw_k<\infty$ and $w_k\ge c_1k^{-s'}$ for some $c_1>0$ and $s'\in\R$, then the Hankel
determinants $D_n$ of $\sum_kw_k\delta_{k^{-p}}$ satisfy
\begin{equation}\label{eq:paperI-logD}
  \log D_n=-pn^2\log n+\kappa_pn^2+o(n^2),\qquad \kappa_p:=\tfrac{3p}2+\log c(p),
\end{equation}
equivalently
\begin{equation}\label{eq:paperI-cesaro}
  \frac2{n^2}\sum_{j=1}^{n-1}(n-j)\log\bigl(j^pb_j\bigr)\longrightarrow\log c(p).
\end{equation}
By Heine's formula, multiplying every weight by a factor in $[e^{-\varepsilon},e^{\varepsilon}]$ changes $D_n$ by a
factor in $[e^{-\varepsilon n},e^{\varepsilon n}]$, so \eqref{eq:paperI-logD} cannot distinguish such weights.
Theorems~\ref{thm:sv} and~\ref{thm:rough} show that the pointwise law does distinguish them: it is stable under
slowly varying perturbations of any size, and unstable under rough ones, however small. We have not found
Theorem~\ref{thm:rough} in the literature.

\emph{Spectral weights.} For a fixed measure with bounded support, Simon \cite[Theorems~3.1 and~3.2]{Simon} relates
the convergence of the Jacobi parameters to weak limits of the measures $p_n^2\diff\mu$, by counting lattice paths
\cite[Proposition~3.3]{Simon}. Our localisation lemma (Lemma~\ref{lem:arcsine}) is the forward half of this argument
at the moving scale $n^{-p}$: the spectral weights of the pure powers, rescaled by $n^p$, converge to the arcsine law
on $[0,4c(p)]$.

\emph{First variation.} Fischer \cite[Lemma~4, (20)--(21)]{Fischer} computes the derivatives of the recurrence
coefficients with respect to a parameter of the orthogonality functional; for a multiplicative change of the weights
these are the formulas of Lemmas~\ref{lem:fv} and~\ref{lem:toda}. What we add in \S\ref{sec:lr} is a bound on the
second derivative with absolute constants, which turns the first-order formulas into estimates for perturbations of
fixed size.

\emph{The Christoffel transformation.} Multiplying $\sigma_{p,s}$ by $x$ gives $\sigma_{p,s+p}$. The Jacobi matrix of
$x\sigma$ is obtained by factorising that of $\sigma$ as $LU$ and reversing the factors \cite[\S1]{DD}, and the
factors form a chain sequence in the sense of Dai, Piao and Qiao \cite[Proposition~2.1]{DPQ}. The second paper
uses this to transfer the law with its rate from $s$ to $s+p$ \cite[Theorem~3.1]{PaperII}. Here the same
factorisation, combined with the product structure of $h_n$, gives the Ces\`aro-level identities of
Section~\ref{sec:christoffel}.

\emph{Special cases.} The moments of $\sigma_{1,s}$ are $\zeta(s+j)$, so Corollary~\ref{cor:p1} concerns Hankel
determinants of zeta values, studied by Monien \cite{Monien} and by Haynes and Zudilin \cite{HaynesZudilin}; their
leading order is \cite[Corollary~1.4]{PaperI}. Our anchor at $p=1$ is the Tricomi--Carlitz measure, whose
orthogonality and recurrence are recorded by Wu, Lin, Xu and Zhao \cite[(1.4)--(1.9)]{WLXZ}; at $p=2$ it is the
Lommel case \cite[Remark~1.6]{PaperII}, which rests on the recurrence recorded by Van Assche \cite[\S3]{VA}.

\subsection{Method}\label{sec:method}
Everything rests on one device. Along a path of weights $w(t)=v\,e^{t\varphi}$, the first variation of $\log h_n$ is
the average of $\varphi$ against the spectral weights $m_n(k)=w_k(t)p_n(x_k)^2$ of the orthonormal polynomial $p_n$
(Lemma~\ref{lem:fv}). Two facts locate these weights: the atoms $k<n/M_0$ are saturated, so that $m_n$ is
exponentially small there (Lemma~\ref{lem:saturation}), and for pure powers $m_n$ lives at indices $k\asymp n$, where its
image under $x\mapsto n^px$ converges to the arcsine law (Lemma~\ref{lem:arcsine}). Each later section is one choice
of $\varphi$. In Section~\ref{sec:sv}, $\varphi=\log(\ell/\ell_1)$ with $\ell_1$ a smoothing of $\ell$; it is small
where $m_n$ lives, which gives stability. In Section~\ref{sec:rough}, $\varphi=\pm\varepsilon$ with the sign of the
second difference $m_n-2m_{n-1}+m_{n-2}$ on a window of indices of size $n$; the linear response then does not
average out, and a second-order bound shows that the remainder cannot cancel it. In Section~\ref{sec:christoffel},
$\varphi=-p\log k$ is multiplication by $x$, which shifts $s$ by $p$. In Section~\ref{sec:s}, $\varphi=-(s_1-s_0)\log
k$ with $s_1-s_0$ small; this gives the derivative in $s$ and a variational two-sided bound for $h_n(s_1)/h_n(s_0)$.

\subsection{Organisation and notation}\label{sec:organisation}
Section~\ref{sec:prelim} collects the tools: conventions, the results imported from \cite{PaperI,PaperII},
slow variation, the first variation and its second-order bounds, and localisation. Sections~\ref{sec:sv},
\ref{sec:rough}, \ref{sec:christoffel} and~\ref{sec:s} prove Theorems~\ref{thm:sv}, \ref{thm:rough},
\ref{thm:s}(i) and~\ref{thm:s}(ii)--(iii) with Corollary~\ref{cor:p1}. Section~\ref{sec:numerics} illustrates the
results numerically, and Section~\ref{sec:open} states open problems and conjectures.

Throughout, $p>0$ is fixed and $c=c(p)$. The atoms are $x_k=k^{-p}$, and $x$ also denotes the sequence $(x_k)$.
Constants $C,C'$ are positive and may change from line to line; the letter $c$ is reserved for $c(p)$, and $c_1$
for the lower-bound constant of the weights in \S\ref{sec:context}. Symbols local to this paper are the smoothing error
$\rho_\ell$ and head margin $M_0$ (Section~\ref{sec:sv} and Lemma~\ref{lem:saturation}), the window constants
$\iota_W,\gamma_W$ and the number $\ia$ (Section~\ref{sec:rough}), the arcsine tail $T_p$ (Lemma~\ref{lem:arcsine}), the exponent $\vartheta_*$ (Lemma~\ref{lem:allM}),
and $\Phi_n,\Theta_n$ (Section~\ref{sec:s}).

% Section 2: First variation and localisation.
\section{First variation and localisation}\label{sec:prelim}

This section collects the tools used in every later section. After fixing conventions (\S\ref{sec:conv}) we restate
the results imported from \cite{PaperI} and \cite{PaperII} (\S\ref{sec:imports}); a reader who accepts these
statements can read the rest of the paper without the earlier ones. We then prove two facts about slowly varying
sequences (\S\ref{sec:prelim-sv}), the first-variation formulas (\S\ref{sec:fv}) with bounds on their second-order
remainders (\S\ref{sec:lr}), and the two localisation properties of the spectral weights (\S\ref{sec:loc}). The proofs
are short and, apart from the second-order bounds, of a standard kind.

\subsection{Conventions}\label{sec:conv}
Throughout, $\sigma=\sum_{k\ge1}w_k\delta_{x_k}$ with distinct atoms $x_k\in(0,1]$, weights $w_k>0$ and
$\|w\|_1:=\sum_kw_k<\infty$, except where a more general positive measure on $[0,\infty)$ is allowed explicitly. Such
a $\sigma$ has infinite support and finite moments, so its monic orthogonal polynomials $P_n$ exist for all $n\ge0$
and satisfy \eqref{eq:rec}. The orthonormal polynomials $p_n:=P_n/\sqrt{h_n}$ satisfy
\begin{equation}\label{eq:rec-orth}
  xp_n=b_{n+1}p_{n+1}+a_np_n+b_np_{n-1}\qquad(n\ge0,\ p_{-1}:=0),
\end{equation}
and therefore $a_n=\int xp_n^2\diff\sigma$ and $b_n=\int xp_{n-1}p_n\diff\sigma$. The Hankel determinants are
$D_n=\prod_{j<n}h_j$. Multiplying all weights by one positive constant changes none of $P_n$, $p_n$, $a_n$, $b_n$.

We regard functions on $\N$ as vectors and write $\langle f,g\rangle:=\sum_kf(k)g(k)$ whenever the sum converges
absolutely, $f(A):=\sum_{k\in A}f(k)$ for $A\subset\N$, and $f|_A:=f\mathbf 1_A$. The \emph{spectral weights} of
$\sigma$ are
\begin{equation}\label{eq:spectral}
  q_n(k):=w_k^{1/2}p_n(x_k),\qquad m_n(k):=q_n(k)^2,\qquad K_n(k):=\sum_{j<n}m_j(k).
\end{equation}
The vectors $q_0,q_1,\dots$ are orthonormal in $\ell^2(\N)$, so each $m_n$ is a probability vector and
$m_n=K_{n+1}-K_n$. By \eqref{eq:rec-orth}, at every atom,
\begin{equation}\label{eq:rec-vec}
  x_kq_n(k)=b_{n+1}q_{n+1}(k)+a_nq_n(k)+b_nq_{n-1}(k),
\end{equation}
so that $a_n=\langle x,m_n\rangle$ and $b_n=\langle xq_{n-1},q_n\rangle\le\frac12\langle x,m_{n-1}+m_n\rangle$. We write
$\Proj_m$ for the orthogonal projection of $\ell^2(\N)$ onto the span of $q_0,\dots,q_{m-1}$; then $K_m(k)$ is the
$k$-th diagonal entry of $\Proj_m$.

\emph{Christoffel's extremal property.} For $y\in\R$ let $\mathsf K_m(y):=\sum_{j<m}p_j(y)^2$. If $\deg P<m$, then
$P=\sum_{j<m}\langle P,p_j\rangle_\sigma p_j$ with $\langle P,p_j\rangle_\sigma:=\int Pp_j\diff\sigma$, and the
Cauchy--Schwarz inequality gives $P(y)^2\le\mathsf K_m(y)\int P^2\diff\sigma$, with equality for
$\langle P,p_j\rangle_\sigma=p_j(y)$. Hence
\begin{equation}\label{eq:christoffel}
  \frac1{\mathsf K_m(y)}=\min\Bigl\{\int P^2\diff\sigma:\ \deg P<m,\ P(y)=1\Bigr\}.
\end{equation}
Every competitor in \eqref{eq:christoffel} with $y=x_k$ has $\int P^2\diff\sigma\ge w_k$, so
$K_m(k)=w_k\mathsf K_m(x_k)\le1$, and in particular $m_n(k)\le1-K_n(k)$.

\subsection{Results imported from the first two papers}\label{sec:imports}
We restate exactly what is used, with the numbering of \cite{PaperI} and \cite{PaperII}. The atoms
$x_k=k^{-p}$ are the case $f(\xi)=\xi^{-p}$, $C=1$ of the atoms $f(k)$ in \cite{PaperI}, since then
$\xi^{p+1}f'(\xi)=-p$.

\begin{theorem}[{\cite[Theorem~1.1 and Corollary~1.2]{PaperI}}]\label{thm:imp-avg}
Let $w_k\ge0$ satisfy $\sum_kw_k<\infty$ and $w_k\ge c_1k^{-s'}$ for all $k\ge1$, with constants $c_1>0$ and
$s'\in\R$. Then the Hankel determinants and the coefficients $b_n$ of $\sum_kw_k\delta_{k^{-p}}$ satisfy
\eqref{eq:paperI-logD} and \eqref{eq:paperI-cesaro}.
\end{theorem}

For $\beta>0$ let $\arc_\beta$ be the arcsine law on $[0,\beta]$, with density $1/(\pi\sqrt{y(\beta-y)})$; it is the
image of $\diff\theta/\pi$ on $(0,\pi)$ under $\theta\mapsto\frac\beta2(1+\cos\theta)$, and
$\arc_\beta([0,y])=\frac2\pi\arcsin\sqrt{y/\beta}$ for $0\le y\le\beta$.

\begin{proposition}[{\cite[Theorem~1.3(i)]{PaperI}}]\label{prop:imp-eq}
Let $\mu^*:=\int_0^1\arc_{4c\tau^{-p}}\diff\tau$, a probability measure on $(0,\infty)$, and let $\mu^*_u$ be its image
under $x\mapsto x^{-1/p}$. Then $\mu^*_u$ has a density $\rho^*$ with $\rho^*=1$ on $(0,\usat]$ and $0<\rho^*<1$ on
$(\usat,\infty)$, where $\usat=\usat(p)=(4c)^{-1/p}$.
\end{proposition}

\begin{definition}[{\cite[Definition~1.1]{PaperII}}]\label{def:analytic}
Let $\vartheta>0$. A sequence $\ell\colon\N\to(0,\infty)$ is \emph{analytic of angle $\vartheta$} if
$\ell(k)=\exp\Lambda(\log k)$ for all $k\ge1$, where, for some $c_0>0$, the function $\Lambda$ is analytic on
$S_\vartheta:=\{\xi\in\mathbb C:\ \Re\xi\ge-c_0,\ |\Im\xi|\le\vartheta\}$, real on $S_\vartheta\cap\R$, bounded on
$S_\vartheta\cap\{\Re\xi\le0\}$, and $\sup\{|\Lambda'(\xi)|:\ \xi\in S_\vartheta,\ \Re\xi\ge T\}\to0$ as $T\to\infty$.
\end{definition}

\begin{theorem}[{\cite[Theorem~1.2]{PaperII}}]\label{thm:imp-analytic}
Let $s>1$, and let $\ell$ be analytic of angle $\vartheta$ for some $\vartheta>0$. Then the coefficients of
$\sigma_{p,s,\ell}$ satisfy $n^pb_n=c(1+O(1/n))$ and $a_n/(2b_n)=1+O(1/n)$; equivalently
$|n^pa_n-2c|+|n^pb_n-c|=O(1/n)$.
\end{theorem}

(The two forms are equivalent: from the first, $n^pa_n=2n^pb_n\cdot a_n/(2b_n)=2c(1+O(1/n))$; from the second,
$a_n/(2b_n)=(2c+O(1/n))/(2c+O(1/n))$.) The case $\ell\equiv1$ is the following.

\begin{theorem}[{\cite[Corollary~1.3]{PaperII}}]\label{thm:imp-pure}
For every $s>1$ the coefficients of $\sigma_{p,s}$ satisfy $n^pb_n=c(1+O(1/n))$ and $a_n/(2b_n)=1+O(1/n)$. In
particular $n^pb_n\to c$ and $n^pa_n\to2c$.
\end{theorem}

Only the last sentence of Theorem~\ref{thm:imp-pure} is used in Sections~\ref{sec:rough}--\ref{sec:s}. The next
statement is quoted only in remarks.

\begin{theorem}[{\cite[Theorem~1.4]{PaperII}}]\label{thm:imp-second}
Let $s>1$.
\begin{enumerate}
\item[\textup{(i)}] If $p\ge1$, then $n^pb_n(s)=c(1-\frac{s-1}{2n}+O(n^{-2}))$ and
  $n^pa_n(s)=2c(1-\frac{s-1+p}{2n}+O(n^{-2}))$.
\item[\textup{(ii)}] If $0<p<1$, then, with $\mathsf k_p:=2p\,c\,\zeta(-p)$,
  \begin{align*}
    n^pb_n(s)&=c\Bigl(1-\frac{s-1}{2n}+\frac{\mathsf k_p}{n^{1+p}}+O\bigl(n^{-1-2p}+n^{-2}\bigr)\Bigr),\\
    n^pa_n(s)&=2c\Bigl(1-\frac{s-1+p}{2n}+\frac{\mathsf k_p}{n^{1+p}}+O\bigl(n^{-1-2p}+n^{-2}\bigr)\Bigr).
  \end{align*}
\end{enumerate}
In both cases $a_n(s)/(2b_n(s))=1-p/(2n)+O(n^{-1-\min(2p,1)})$.
\end{theorem}

The Christoffel transformation is the subject of \cite[Section~3]{PaperII}. Let $\sigma$ be a positive Borel measure
on $[0,\infty)$ with infinite support and finite moments. By \cite[Lemma~3.2]{PaperII}, $P_n$ has $n$ simple zeros in
$(0,\infty)$, so $(-1)^nP_n(0)>0$, and one may define
\begin{equation}\label{eq:chain}
  u_n:=-\frac{P_{n+1}(0)}{P_n(0)}>0\quad(n\ge0),\qquad l_0:=0,\qquad l_n:=\frac{b_n^2}{u_{n-1}}>0\quad(n\ge1).
\end{equation}

\begin{lemma}[{\cite[Lemmas~3.3 and~3.4]{PaperII}}]\label{lem:imp-chain}
Let $\sigma$ be as above, with coefficients $a_n,b_n$, and let $\tilde a_n,\tilde b_n$ be the coefficients of
$x\sigma$, $\diff(x\sigma):=x\diff\sigma(x)$.
\begin{enumerate}
\item[\textup{(i)}] $a_n=u_n+l_n$ for $n\ge0$, $b_n^2=u_{n-1}l_n$ for $n\ge1$; $x\sigma$ has infinite support, its
  monic orthogonal polynomials are $(P_{n+1}+u_nP_n)/x$, and
  \[
    \tilde a_n=u_n+l_{n+1}\quad(n\ge0),\qquad\tilde b_n^2=u_nl_n\quad(n\ge1).
  \]
\item[\textup{(ii)}] If $c'>0$, $n^pa_n\to2c'$ and $n^pb_n\to c'$, then $n^pu_n\to c'$ and $n^pl_n\to c'$.
\end{enumerate}
\end{lemma}

\begin{theorem}[{\cite[Theorem~3.1]{PaperII}}, without rate]\label{thm:imp-shift}
Let $\sigma$ be as above and $c'>0$. If $n^pa_n(\sigma)\to2c'$ and $n^pb_n(\sigma)\to c'$, then the same holds for
$x\sigma$. For the atoms $k^{-p}$ one has $x\,\sigma_{p,s,\ell}=\sigma_{p,s+p,\ell}$.
\end{theorem}

\begin{proposition}[{\cite[Remark~1.6]{PaperII}}]\label{prop:imp-lommel}
For $n\ge1$ the coefficients of $\sigma_{2,2}$ are $b_n=\pi^2(4n+1)^{-1}((4n-1)(4n+3))^{-1/2}$ and
$a_n=2\pi^2(4n+1)^{-1}(4n+5)^{-1}$.
\end{proposition}

The last import is a bound for the spectral weights of the pure powers at indices beyond $Mn$, that is, near the
accumulation point $x=0$ of the atoms. It is used only in Section~\ref{sec:s}. Recall the set $\Sp$ of
\eqref{eq:Sp}.

\begin{proposition}[\tailbound]\label{prop:imp-tail}
Let $s\in\Sp$ \textup{(}that is, $1<s<1+2p$, or $p>1$ and $s>1$\textup{)}, and let $m_n$ be the spectral weights of
$\sigma_{p,s}$. There are $C,\gamma_*>0$ and $n_0$ such that $m_{n'}((Mn,\infty))\le CM^{-p/2}$ for all $n\ge n_0$,
$n'\in\{n-1,n\}$ and $1\le M\le\gamma_*n^{2/p}$.
\end{proposition}

For $p>1$, \tailbound{} also gives a bound for all $M\ge1$ with any exponent below $\min(p/2,s-1)$. We do not need it:
Lemma~\ref{lem:allM} below derives a bound for all $M\ge1$ from Proposition~\ref{prop:imp-tail} alone, with a smaller
exponent, and this suffices for Section~\ref{sec:s}.

\subsection{Slow variation}\label{sec:prelim-sv}
Let $\ell$ be slowly varying and put $\osc^*_n:=\sup_{m\ge n}\osc_m(2)$. Each $\osc_m(2)$ is a maximum of finitely
many numbers and $\osc_m(2)\to0$, so $\osc^*_n$ is finite, non-increasing, and tends to $0$.

\begin{lemma}[dyadic chaining and Potter bounds]\label{lem:potter}
Let $\ell$ be slowly varying.
\begin{enumerate}
\item[\textup{(i)}] For all $n\ge1$ and all integers $k\ge n$,
  $|\log\ell(k)-\log\ell(n)|\le\osc^*_n\bigl(1+\log_2(k/n)\bigr)$.
\item[\textup{(ii)}] For every $\eta>0$ there is $A_\eta<\infty$ with $e^{-A_\eta}k^{-\eta}\le\ell(k)\le
  e^{A_\eta}k^{\eta}$ for all $k\ge1$.
\end{enumerate}
\end{lemma}

\begin{proof}
(i) Let $j:=\lfloor\log_2(k/n)\rfloor$, so that $2^jn\le k<2^{j+1}n$. For $0\le i<j$ the integer $2^{i+1}n$ lies in
$[2^in/2,2\cdot2^in]$, so $|\log\ell(2^{i+1}n)-\log\ell(2^in)|\le\osc_{2^in}(2)\le\osc^*_n$; also
$|\log\ell(k)-\log\ell(2^jn)|\le\osc^*_n$. Summing these $j+1$ steps gives the claim, since $j\le\log_2(k/n)$.

(ii) Choose $m$ with $\osc^*_m\le\eta\log2$. For $k\ge m$, part (i) gives $|\log\ell(k)|\le|\log\ell(m)|+\osc^*_m+
\eta\log k$. Take $A_\eta:=\max(|\log\ell(m)|+\osc^*_m,\max_{k<m}|\log\ell(k)|)$.
\end{proof}

\begin{remark}\label{rem:lambda2}
If $\osc_n(2)\to0$, then $\osc_n(\lambda)\to0$ for every $\lambda\ge1$: Lemma~\ref{lem:potter}(i) uses only
$\osc(2)$, and for $n\le k\le\lambda n$ it gives $|\log\ell(k)-\log\ell(n)|\le\osc^*_n(1+\log_2\lambda)$, while for
$n/\lambda\le k<n$, applied with base point $k$, it gives $|\log\ell(n)-\log\ell(k)|\le\osc^*_{\lceil
n/\lambda\rceil}(1+\log_2\lambda)$.
\end{remark}

By Lemma~\ref{lem:potter}(ii) with $\eta<s-1$, $\sum_kk^{-s}\ell(k)<\infty$ for every $s>1$, so $\sigma_{p,s,\ell}$
is a measure of the kind fixed in \S\ref{sec:conv}.

\subsection{First variation}\label{sec:fv}
Fix atoms and weights $v=(v_k)$ as in \S\ref{sec:conv}, a bounded real sequence $\varphi$ on $\N$, and put
\[
  w(t):=v\,e^{t\varphi}\qquad(t\in\R).
\]
All quantities of \S\ref{sec:conv} computed for $w(t)$ carry the argument $t$ when needed, and a dot denotes
$\diff/\diff t$. Formulas of this kind appear in Fischer \cite[Lemma~4]{Fischer}; we include short proofs in the form
needed here. For $\varphi=x$ they reduce to the Toda equations.

\emph{Regularity.} The moments $\sum_kv_ke^{t\varphi(k)}x_k^i$ are $C^\infty$ in $t$, with derivatives computed term
by term, since the $r$-th termwise derivative is bounded by $v_k\|\varphi\|_\infty^re^{|t|\|\varphi\|_\infty}$. The
coefficients of $P_n$ solve a linear system whose determinant is $D_n>0$, so they are $C^\infty$, and so are $h_n$,
$a_n$, $b_n^2$. The vectors $w(t)^{1/2}x^j$ are $C^\infty$ maps into $\ell^2(\N)$ (the $r$-th derivative is
$(\varphi/2)^rw(t)^{1/2}x^j$), and $q_0,\dots,q_n$ arise from finitely many of them by the Gram--Schmidt process, so
they are $C^\infty$ maps into $\ell^2(\N)$ as well.

\begin{lemma}[first variation]\label{lem:fv}
For $n\ge0$,
\begin{align}
  \frac{\diff}{\diff t}\log h_n&=\langle\varphi,m_n\rangle, \label{eq:fv-h}\\
  \frac{\diff}{\diff t}\log b_n^2&=\langle\varphi,m_n-m_{n-1}\rangle\qquad(n\ge1), \label{eq:fv-b}\\
  \frac{\diff}{\diff t}a_n&=\langle(x-a_n)\varphi,m_n\rangle-2b_n\langle\varphi,q_{n-1}q_n\rangle. \label{eq:fv-a}
\end{align}
The right-hand sides do not change if a constant is added to $\varphi$.
\end{lemma}

\begin{proof}
For a polynomial $g(x;t)$ whose coefficients are $C^1$ in $t$,
$\frac{\diff}{\diff t}\sum_kw_k(t)g(x_k;t)=\sum_kw_k(t)\partial_tg(x_k;t)+\sum_k\varphi(k)w_k(t)g(x_k;t)$, because the
left side is a finite combination of moments. Apply this to $g=P_n^2$: since $P_n$ is monic, $\partial_tP_n$ has
degree less than $n$ and is orthogonal to $P_n$, so $\dot h_n=\sum_k\varphi(k)w_kP_n(x_k)^2=h_n\langle\varphi,m_n\rangle$,
which is \eqref{eq:fv-h}; \eqref{eq:fv-b} follows from $b_n^2=h_n/h_{n-1}$. For \eqref{eq:fv-a}, expand
$\dot p_n=\sum_{j\le n}\mathsf c_jp_j$. Differentiating $\int p_np_j\diff\sigma=\delta_{nj}$ gives
$\mathsf c_j=-\langle\varphi,q_nq_j\rangle$ for $j<n$ (the term $\int p_n\dot p_j\diff\sigma$ vanishes because
$\deg\dot p_j\le j<n$) and $\mathsf c_n=-\frac12\langle\varphi,m_n\rangle$. Then
\[
  \dot a_n=2\int xp_n\dot p_n\diff\sigma+\langle x\varphi,m_n\rangle=2\mathsf c_na_n+2\mathsf c_{n-1}b_n
  +\langle x\varphi,m_n\rangle,
\]
by \eqref{eq:rec-orth}, which is \eqref{eq:fv-a}. Finally $\langle1,m_n-m_{n-1}\rangle=0$,
$\langle x-a_n,m_n\rangle=0$ and $\langle q_{n-1},q_n\rangle=0$.
\end{proof}

\begin{lemma}[Toda form]\label{lem:toda}
Put $\cE_n:=b_{n+1}q_nq_{n+1}$ \textup{(}a vector on $\N$; $\cE_{-1}:=0$\textup{)} and $\Delta\cE_n:=\cE_n-\cE_{n-1}$. Then
at every atom
\begin{equation}\label{eq:toda-id}
  \Delta\cE_n=(x-a_n)m_n-2b_nq_{n-1}q_n=2\cE_n-(x-a_n)m_n,
\end{equation}
and $\dot a_n=\langle\varphi,\Delta\cE_n\rangle$. Moreover $|\Delta\cE_n|\le\frac12\max(b_n,b_{n+1})
(m_{n-1}+2m_n+m_{n+1})$ at every atom.
\end{lemma}

\begin{proof}
Multiplying \eqref{eq:rec-vec} by $q_n(k)$ gives $(x_k-a_n)m_n(k)=\cE_n(k)+\cE_{n-1}(k)$, which yields both forms of
\eqref{eq:toda-id}; pairing the first with $\varphi$ and comparing with \eqref{eq:fv-a} gives
$\dot a_n=\langle\varphi,\Delta\cE_n\rangle$. The bound follows from $|q_jq_{j+1}|\le\frac12(m_j+m_{j+1})$.
\end{proof}

\begin{lemma}[derivative of the spectral vectors]\label{lem:qdot}
For $m\ge0$, with $\alpha_m:=\langle\varphi,m_m\rangle$,
\[
  \dot q_m=\tfrac12\varphi q_m-\Proj_m(\varphi q_m)-\tfrac12\alpha_mq_m .
\]
\end{lemma}

\begin{proof}
Since $q_m=w(t)^{1/2}p_m(x)$ and $\deg\dot p_m\le m$, we have $\dot q_m=\frac12\varphi q_m+r$ with $r$ in the span of
$q_0,\dots,q_m$. For $j<m$ the same reasoning gives $\dot q_j-\frac12\varphi q_j\perp q_m$, so differentiating
$\langle q_m,q_j\rangle=0$ gives $\langle r,q_j\rangle=-\langle\varphi q_m,q_j\rangle$. Differentiating
$\langle q_m,q_m\rangle=1$ gives $\langle r,q_m\rangle=-\frac12\alpha_m$. Hence $r=-\Proj_m(\varphi q_m)-\frac12\alpha_mq_m$.
\end{proof}

\subsection{Second-order bounds}\label{sec:lr}
The first-variation formulas are exact derivatives; to use them for a perturbation $e^\varphi$ of fixed size one
needs the second derivative. The next lemma bounds it with absolute constants, for every base measure. It is the
reason why the constructions of Section~\ref{sec:rough} need no information about the perturbed measures.

For a bounded real $\varphi$ put $\varepsilon:=\frac12(\sup\varphi-\inf\varphi)$. Write
$\Delta^2m_N:=m_N-2m_{N-1}+m_{N-2}$ and $F_N:=\log(b_N^2/b_{N-1}^2)$ for $N\ge2$. A superscript $v$ denotes quantities
of the base weights $v$, and $w:=w(1)=ve^\varphi$.

\begin{lemma}[linear response with quadratic remainder]\label{lem:lr}
In the setting of \S\ref{sec:fv}, with $\varepsilon=\frac12\osc\varphi$:
\begin{enumerate}
\item[\textup{(i)}] $|\log b_n^2(w)-\log b_n^2(v)-\langle\varphi,m_n^v-m_{n-1}^v\rangle|\le\varepsilon^2$ for $n\ge1$;
\item[\textup{(ii)}] $|F_N(w)-F_N(v)-\langle\varphi,\Delta^2m_N^v\rangle|\le2\varepsilon^2$ for $N\ge2$;
\item[\textup{(iii)}] $|a_n(w)-a_n(v)-\langle\varphi,\Delta\cE_n^v\rangle|\le(e^\varepsilon-1-\varepsilon)
  \bigl(b_n(v)+b_{n+1}(v)\bigr)$ for $n\ge0$.
\end{enumerate}
\end{lemma}

\begin{proof}
\emph{Reduction.} With $\bar c:=\frac12(\sup\varphi+\inf\varphi)$ we have $ve^{t\varphi}=e^{t\bar c}ve^{t(\varphi-\bar
c)}$, and constant factors change none of the quantities involved; the linear terms do not change either
(Lemma~\ref{lem:fv}). So we may assume $|\varphi|\le\varepsilon$. Since $w(t+h)=w(t)e^{h\varphi}$ and $w(t)$ is again
admissible, every bound below proved at $t=0$ for an arbitrary base holds at every $t$. We use Taylor's formula
$g(1)-g(0)-\dot g(0)=\int_0^1(1-t)\ddot g(t)\diff t$.

\emph{Second derivative of $\log h_m$.} By Lemma~\ref{lem:fv} and Lemma~\ref{lem:qdot},
\[
  \dot\alpha_m=2\langle\varphi q_m,\dot q_m\rangle=\|\varphi q_m\|^2-\alpha_m^2-2\|\Proj_m(\varphi q_m)\|^2,
\]
since $\Proj_m$ is an orthogonal projection. Because $q_m$ is orthogonal to the range of $\Proj_m$, the vectors
$\Proj_m(\varphi q_m)$, $\alpha_mq_m$ and the rest of $\varphi q_m$ are pairwise orthogonal, and Bessel's inequality
gives $\|\Proj_m(\varphi q_m)\|^2\le\|\varphi q_m\|^2-\alpha_m^2=\Var_{m_m}(\varphi)$. Hence
\[
  |\dot\alpha_m|\le\Var_{m_m}(\varphi)\le\sum_k\varphi(k)^2m_m(k)\le\varepsilon^2 .
\]

\emph{(i) and (ii).} By \eqref{eq:fv-b}, $\frac{\diff^2}{\diff t^2}\log b_n^2=\dot\alpha_n-\dot\alpha_{n-1}$, of absolute
value at most $2\varepsilon^2$, and $\ddot F_N=\dot\alpha_N-2\dot\alpha_{N-1}+\dot\alpha_{N-2}$, of absolute value at most
$4\varepsilon^2$. Taylor's formula gives (i) and (ii), since $\int_0^1(1-t)\diff t=\frac12$.

\emph{(iii).} Put $\gamma_m:=\langle\varphi q_m,q_{m+1}\rangle$. By \eqref{eq:rec-vec},
$\langle(x-a_n)\varphi,m_n\rangle=b_{n+1}\gamma_n+b_n\gamma_{n-1}$, so \eqref{eq:fv-a} reads
$\dot a_n=b_{n+1}\gamma_n-b_n\gamma_{n-1}$. By \eqref{eq:fv-b}, $\dot b_m=\frac12(\alpha_m-\alpha_{m-1})b_m$; in particular
$|\dot b_m|\le\varepsilon b_m$, so $b_m(t)\le e^{\varepsilon t}b_m(0)$. Put $u:=\varphi q_m$, $u':=\varphi q_{m+1}$,
$\Proj:=I-\Proj_{m+2}$, $S_m:=\langle\Proj_mu,\Proj_mu'\rangle$ and $R_m:=\langle\Proj u,\Proj u'\rangle$. Decomposing
along the range of $\Proj_m$, $q_m$, $q_{m+1}$ and the range of $\Proj$ (pairwise orthogonal), and using
Lemma~\ref{lem:qdot} twice, one finds
\[
  \langle\dot q_m,u'\rangle=\tfrac12\langle u,u'\rangle-S_m-\tfrac12\alpha_m\gamma_m,\qquad
  \langle u,\dot q_{m+1}\rangle=\tfrac12\langle u,u'\rangle-S_m-\alpha_m\gamma_m-\tfrac12\alpha_{m+1}\gamma_m,
\]
and $\langle u,u'\rangle=S_m+\alpha_m\gamma_m+\gamma_m\alpha_{m+1}+R_m$. Adding,
$\dot\gamma_m=R_m-S_m+\frac12(\alpha_{m+1}-\alpha_m)\gamma_m$, and therefore
\[
  \ddot a_n=b_{n+1}\Gamma_n-b_n\Gamma_{n-1},\qquad \Gamma_m:=R_m-S_m+(\alpha_{m+1}-\alpha_m)\gamma_m .
\]
We claim $|\Gamma_m|\le\varepsilon^2$. By the Cauchy--Schwarz inequality, first in $\ell^2$ and then in $\R^2$,
$|R_m-S_m|\le(\|\Proj u\|^2+\|\Proj_mu\|^2)^{1/2}(\|\Proj u'\|^2+\|\Proj_mu'\|^2)^{1/2}$. By the orthogonal decomposition,
$\|\Proj u\|^2+\|\Proj_mu\|^2=\|u\|^2-\alpha_m^2-\gamma_m^2\le\varepsilon^2-\alpha_m^2-\gamma_m^2$, and likewise for $u'$
with $\alpha_{m+1}$. By the inequality between the geometric and arithmetic means,
$|R_m-S_m|\le\varepsilon^2-\gamma_m^2-\frac12(\alpha_m^2+\alpha_{m+1}^2)$, while
$|(\alpha_{m+1}-\alpha_m)\gamma_m|\le\gamma_m^2+\frac14(|\alpha_m|+|\alpha_{m+1}|)^2\le\gamma_m^2+\frac12(\alpha_m^2+\alpha_{m+1}^2)$.
Adding proves the claim. Hence $|\ddot a_n(t)|\le\varepsilon^2(b_n(t)+b_{n+1}(t))\le\varepsilon^2e^{\varepsilon t}(b_n+b_{n+1})(v)$,
and Taylor's formula with $\dot a_n(0)=\langle\varphi,\Delta\cE_n^v\rangle$ (Lemma~\ref{lem:toda}) gives (iii), because
$\int_0^1(1-t)\varepsilon^2e^{\varepsilon t}\diff t=e^\varepsilon-1-\varepsilon$.
\end{proof}

\begin{remark}\label{rem:lr-sharp}
The constants in Lemma~\ref{lem:lr} do not depend on $n$, on the atoms, or on $v$. The pointwise bound
$|\ddot a_n|\le\varepsilon^2(b_n+b_{n+1})$ is attained: for two atoms $\pm1$, $n=0$, weights proportional to
$(e^\theta,e^{-\theta})$ and $\varphi=\pm\varepsilon$ on the two atoms, $a_0=\tanh\theta$ and $b_1=1/\cosh\theta$, and
$|\ddot a_0|=2\varepsilon^2|\tanh\theta|/\cosh^2\theta$ equals $\varepsilon^2(b_0+b_1)$ when $\sinh\theta=\pm1$. (This
example has atoms $\pm1$; the proof of Lemma~\ref{lem:lr} uses only that the atoms are bounded.)
\end{remark}

\subsection{Localisation}\label{sec:loc}
Two facts locate the spectral weights. The first holds for all weights and says that the atoms $x_k$ with $k$ much
smaller than $n$ are almost fully occupied by $p_0,\dots,p_{n-1}$. For $p>0$ and $0<y<1$ let
\begin{equation}\label{eq:gp}
  g_p(y):=(p-y)\log\frac1y+(1-y)\log(1-y)-p+(p-1)y,\qquad
  g_p^+:=\begin{cases}g_p,&p\ge1,\\ g_p-\log\frac1p,&0<p<1.\end{cases}
\end{equation}
For $y<p/2$ the first term is at least $\frac p2\log\frac1y$ and the others are bounded, so $g_p^+(y)\to\infty$ as
$y\to0^+$. The \emph{head margin} $M_0=M_0(p)$ is the least integer $M_0\ge4$ such that $g_p^+\ge1$ on $(0,2/M_0]$.

\begin{lemma}[saturation of the head]\label{lem:saturation}
Let $\sigma=\sum_kw_k\delta_{k^{-p}}$ with $w_k>0$ and $\|w\|_1<\infty$. For integers $1\le k<n$ and $y:=k/n$,
\[
  m_n(k)\le1-K_n(k)\le\frac{\|w\|_1}{w_k}\Bigl(\frac{n+1}k\Bigr)^{2(1-p)_+}e^{-2ng_p^+(y)} .
\]
\end{lemma}

\begin{proof}
The first inequality was noted after \eqref{eq:christoffel}. Use the test polynomial
$P(x):=\prod_{j\le n,\,j\ne k}(x-x_j)/(x_k-x_j)$ in \eqref{eq:christoffel}: it has degree $n-1$, $P(x_k)=1$, and it
vanishes at $x_j$ for $j\le n$, $j\ne k$. For $i>n$ and $j\le n$ we have $0<x_i<x_j$, so $|P(x_i)|\le\mathcal
P:=\prod_{j\le n,\,j\ne k}x_j/|x_k-x_j|$. Hence $1/\mathsf K_n(x_k)\le w_k+\|w\|_1\mathcal P^2$, and
therefore $1-K_n(k)\le\|w\|_1\mathcal P^2/w_k$. To bound $\mathcal P$, note $x_j/|x_k-x_j|=k^p/|k^p-j^p|$ and
$|k^p-j^p|\ge\varsigma_p\max(j,k)^{p-1}|k-j|$ with $\varsigma_p:=\min(1,p)$ (directly for $p\ge1$, by the mean value
theorem for $p<1$). Thus
\[
  \mathcal P\le\varsigma_p^{-(n-1)}\frac{k^{k-1}}{(k-1)!}\cdot\frac{k^{p(n-k)}(n!/k!)^{1-p}}{(n-k)!}.
\]
Use $k^{k-1}/(k-1)!\le e^k$, $\log(n-k)!\ge(n-k)\log(n-k)-(n-k)$, and $I_{k,n}\le\log(n!/k!)\le
I_{k,n}+\log\frac{n+1}k$ with $I_{k,n}:=\int_k^n\log t\diff t$, taking the lower bound when $p\ge1$ and the upper bound
when $p<1$. With $k=yn$ the terms in $n\log n$ cancel, and
$k+p(n-k)\log k+(1-p)I_{k,n}-(n-k)\log(n-k)+(n-k)=-ng_p(y)$. Hence $\log\mathcal P\le-ng_p(y)+(n-1)\log(1/\varsigma_p)
+(1-p)_+\log\frac{n+1}k$, and squaring gives the claim.
\end{proof}

\begin{corollary}[head mass]\label{cor:head}
Let $w$ be as in Lemma~\ref{lem:saturation} with $w_k\ge c_1k^{-A}$ for all $k$, where $c_1>0$ and $A\ge0$. For every integer $m\ge2$,
\[
  m_m\bigl([1,2m/M_0]\bigr)\le\frac{\|w\|_1}{c_1}\,m^{A+1}(m+1)^2e^{-2m}.
\]
\end{corollary}

\begin{proof}
For $1\le k\le2m/M_0$ we have $k\le m/2<m$ and $g_p^+(k/m)\ge1$, and $((m+1)/k)^{2(1-p)_+}\le(m+1)^2$. Apply
Lemma~\ref{lem:saturation} and sum over fewer than $m$ indices.
\end{proof}

The second fact concerns the pure powers $\sigma_{p,s}$ and is where Theorem~\ref{thm:imp-pure} enters. Put
\begin{equation}\label{eq:Tp}
  T_p(M):=\arc_{4c}\bigl([0,M^{-p})\bigr)=\frac2\pi\arcsin\sqrt{\min\bigl(1,M^{-p}/(4c)\bigr)}\qquad(M\ge1),
\end{equation}
which decreases to $0$ as $M\to\infty$.

\begin{lemma}[arcsine localisation]\label{lem:arcsine}
Let $s>1$, let $m_n$ be the spectral weights of $\sigma_{p,s}$, and fix $i\in\mathbb Z$. Let $\nu_{n,i}$ be the image
of $m_{n+i}$ under $k\mapsto n^px_k$, a probability measure on $(0,\infty)$. Then $\nu_{n,i}\to\arc_{4c}$ weakly as
$n\to\infty$. In particular $m_{n+i}((Mn,\infty))\to T_p(M)$ for every $M\ge1$.
\end{lemma}

\begin{proof}
Let $\cJ$ be the Jacobi matrix of $\sigma_{p,s}$, with diagonal $(a_0,a_1,\dots)$ and off-diagonal
$(b_1,b_2,\dots)$. By \eqref{eq:rec-vec}, $\int y^j\diff\nu_{n,i}(y)=n^{pj}\langle x^jq_{n+i},q_{n+i}\rangle
=n^{pj}(\cJ^j)_{n+i,n+i}$, which is a sum over lattice paths of length $j$ from $n+i$ to itself, with steps $0,\pm1$,
of products of entries $a_m,b_m$ with $|m-n|\le j+|i|+1$; compare \cite[Proposition~3.3]{Simon}. For
$n>j+|i|+1$ the set of paths does not depend on $n$, and by Theorem~\ref{thm:imp-pure} $n^pa_m\to2c$ and $n^pb_m\to
c$ for each such $m$, since $(m/n)^p\to1$. So the moment converges to the corresponding sum for the doubly infinite
matrix $J_0$ with diagonal $2c$ and off-diagonal $c$, that is, to the constant Fourier coefficient of
$(c(2+e^{i\theta}+e^{-i\theta}))^j$, which equals $\frac1\pi\int_0^\pi(2c(1+\cos\theta))^j\diff\theta=\int
y^j\diff\arc_{4c}$.

Moment convergence implies weak convergence here, because the limit has compact support. The second moments
are bounded, so $(\nu_{n,i})_n$ is tight; let $\nu'$ be a subsequential weak limit. Since all moments of $\nu_{n,i}$ are
bounded, each $y^j$ is uniformly integrable, so $\int y^j\diff\nu'=\int y^j\diff\arc_{4c}\le(4c)^j$. Then
$\nu'((4c+\delta,\infty))\le(4c/(4c+\delta))^{2j}\to0$ for every $\delta>0$, so $\nu'$ lives on $[0,4c]$, where
polynomials are dense in the continuous functions; hence $\nu'=\arc_{4c}$. Finally,
$m_{n+i}((Mn,\infty))=\nu_{n,i}([0,M^{-p}))$, and $\arc_{4c}$ has no atoms.
\end{proof}

Beyond the range of Proposition~\ref{prop:imp-tail}, a soft bound suffices. It uses two elementary facts about a
real polynomial $P$ of degree at most $n$, with $\|P\|:=\max_{[-1,1]}|P|$ and $T_j$ the Chebyshev polynomials,
$T_j(\cos\theta)=\cos j\theta$:
\begin{equation}\label{eq:cheb}
  \max_{[-1,1]}|P'|\le2n^3\|P\|,\qquad |P(y)|\le2(n+1)\,|T_n(y)|\,\|P\|\quad(|y|\ge1).
\end{equation}
Indeed, $P=\sum_{j\le n}c_jT_j$ with $c_j=\frac{2-\delta_{j0}}\pi\int_0^\pi P(\cos\theta)\cos j\theta\diff\theta$, so
$|c_j|\le2\|P\|$. On $[-1,1]$, $|T_j'(\cos\theta)|=|j\sin j\theta/\sin\theta|\le j^2$, since $|\sin j\theta|\le j|\sin\theta|$; this
gives the first bound. For $|y|\ge1$, $|T_j(y)|=\cosh(j\operatorname{arccosh}|y|)$ is non-decreasing in $j$, which gives
the second.

\begin{lemma}[polynomial tail localisation]\label{lem:TL}
Let $\sigma=\sum_kw_k\delta_{k^{-p}}$ with $c'k^{-s}\le w_k\le C'k^{-s}$ for all $k$, where $s>1$ and $0<c'\le C'$. For
$n\ge2$ put $T':=\lceil16pn^3\rceil$ and $T'':=\lceil T'n^{2/p}\rceil$. Then $K_{n+1}(k)\le68^2(n+1)^2(C'/c')(T''/k)^s$ for
every $k\ge T''$, and for every integer $T\ge T''$ and $n'\in\{n-1,n\}$,
\[
  m_{n'}\bigl((T,\infty)\bigr)\le\frac{68^2(n+1)^2(C'/c')}{s-1}\,T''^{\,s}\,T^{1-s}.
\]
\end{lemma}

\begin{proof}
Put $\bar x:=x_{T'}$ and $\underline x:=x_{T''}$, so that $\underline x/\bar x\le(T'/T'')^p\le n^{-2}\le\frac14$. Let $\deg P\le n$ and write $\|P\|_{[\underline x,\bar x]}$
for the maximum of $|P|$ on $[\underline x,\bar x]$. Mapping $[\underline x,\bar x]$ affinely onto $[-1,1]$, \eqref{eq:cheb} gives $|P'|\le
4n^3\|P\|_{[\underline x,\bar x]}/(\bar x-\underline x)\le8n^3\|P\|_{[\underline x,\bar x]}/\bar x$ on $[\underline x,\bar x]$. Consecutive atoms in $[\underline x,\bar x]$ are at distance $x_j-x_{j+1}\le
pj^{-p-1}\le p\bar x/T'$, so every point of $[\underline x,\bar x]$ is within $p\bar x/T'$ of an atom, and
$\|P\|_{[\underline x,\bar x]}\le\max|P(x_j)|+(8pn^3/T')\|P\|_{[\underline x,\bar x]}$, the maximum over $T'\le j\le T''$. Since $T'\ge16pn^3$,
$\|P\|_{[\underline x,\bar x]}\le2\max_{T'\le j\le T''}|P(x_j)|$. For $0\le x\le\underline x$ the affine image $y=(\underline x+\bar x-2x)/(\bar x-\underline x)$ satisfies
$1\le y\le1+4\underline x/\bar x\le1+4n^{-2}$, and $T_n(1+\eta)=\cosh(n\operatorname{arccosh}(1+\eta))\le e^{n\sqrt{2\eta}}\le e^{\sqrt8}<17$ for
$\eta\le4n^{-2}$. By \eqref{eq:cheb}, $|P(x)|\le34(n+1)\|P\|_{[\underline x,\bar x]}$. Hence, if $k\ge T''$ and $P(x_k)=1$, then
$1\le68(n+1)\max_{T'\le j\le T''}|P(x_j)|$, so $\int P^2\diff\sigma\ge\min_{j\le T''}w_j/(68(n+1))^2\ge c'T''^{-s}/(68(n+1))^2$.
By \eqref{eq:christoffel}, $K_{n+1}(k)=w_k\mathsf K_{n+1}(x_k)\le68^2(n+1)^2(C'/c')(T''/k)^s$. Finally
$m_{n'}\le K_{n'+1}\le K_{n+1}$ for $n'\le n$, and $\sum_{k>T}k^{-s}\le T^{1-s}/(s-1)$.
\end{proof}

\begin{lemma}[one power for all $M$]\label{lem:allM}
Let $s\in\Sp$, let $m_n$ be the spectral weights of $\sigma_{p,s}$, and put
\[
  \vartheta_*=\vartheta_*(p,s):=\frac{s-1}{4+(2+2/p)s}.
\]
Then $0<\vartheta_*<\min(p/2,s-1)$, $\vartheta_*$ increases with $s$, and there are $C$ and $n_0$ such that
$m_{n'}((Mn,\infty))\le CM^{-\vartheta_*}$ for all $n\ge n_0$, $n'\in\{n-1,n\}$ and $M\ge1$.
\end{lemma}

\begin{proof}
Put $B:=3+(2+2/p)s$, so that $\vartheta_*=(s-1)/(B+1)$, and $\gamma_1:=1/\vartheta_*=(B+1)/(s-1)$; then
$\gamma_1>2+2/p$, and $\vartheta_*<(s-1)p/(2s)<p/2$. The map $M\mapsto m_{n'}((Mn,\infty))$ is non-increasing. For $1\le
M\le\gamma_*n^{2/p}$ use Proposition~\ref{prop:imp-tail} and $M^{-p/2}\le M^{-\vartheta_*}$. For $\gamma_*n^{2/p}\le M\le
n^{\gamma_1}$, monotonicity and Proposition~\ref{prop:imp-tail} give the bound $C\gamma_*^{-p/2}n^{-1}\le C\gamma_*^{-p/2}M^{-\vartheta_*}$,
since $M^{\vartheta_*}\le n$. For $M\ge n^{\gamma_1}$ apply Lemma~\ref{lem:TL} with $c'=C'=1$ and $T:=\lfloor Mn\rfloor\ge
Mn/2$; since $T''\le(16p+2)n^{3+2/p}$ and $\gamma_1+1>3+2/p$, we have $T\ge T''$ for large $n$, and
$\{k>Mn\}\subset\{k>T\}$. This gives $m_{n'}((Mn,\infty))\le Cn^2n^{(3+2/p)s}(Mn)^{1-s}=Cn^BM^{-(s-1)}$, and
$n^BM^{-(s-1)}=n^BM^{-B\vartheta_*}M^{-\vartheta_*}\le M^{-\vartheta_*}$ because $M^{\vartheta_*}\ge n$.
\end{proof}

\begin{lemma}\label{lem:usat}
$\usat(p)<1$ for every $p>0$.
\end{lemma}

\begin{proof}
By Proposition~\ref{prop:imp-eq}, $\mu^*_u((0,\usat])=\usat$. In the variable $x=u^{-p}$ the set $u\le\usat$ is
$x\ge4c$, so $\usat=1-\mu^*((0,4c))$. For $\tau\in[\frac12,1]$ the law $\arc_{4c\tau^{-p}}$ gives the mass
$\frac2\pi\arcsin(\tau^{p/2})\ge\frac2\pi\arcsin(2^{-p/2})$ to $[0,4c)$, so $\mu^*((0,4c))>0$.
\end{proof}

% Section 3: every slowly varying factor (Theorem 1.2).
\section{Every slowly varying factor}\label{sec:sv}

This section proves Theorem~\ref{thm:sv}. A slowly varying sequence is close to an analytic one: Gaussian smoothing
of $\log\ell$ in the logarithmic variable produces an entire function $\Lambda_\ell$ whose derivatives tend to $0$ in
every horizontal strip and which differs from $\log\ell$ by $o(1)$ at the integers (\S\ref{sec:smoothing}). The factor
$\ell_1:=e^{\Lambda_\ell(\log\cdot)}$ is therefore analytic of every angle, and Theorem~\ref{thm:imp-analytic} applies to
it. It remains to compare the weights $k^{-s}\ell(k)$ and $k^{-s}\ell_1(k)$, whose ratio tends to $1$, along the path
of \S\ref{sec:fv} with $\varphi=\log(\ell/\ell_1)$ (\S\ref{sec:stability}). The first variation is an average of
$\varphi$ against spectral weights that live at indices $k\ge n/M_0$, up to an exponentially small head, and there
$|\varphi|$ is at most the smoothing error.

One cannot compare $\ell$ directly with the constant $1$: $\ell$ may be unbounded and need not be close to $1$
anywhere, and controlling $\log\ell(k)-\log\ell(n)$ against the spectral weights far beyond $k\asymp n$ would require
information about the accumulation point of the atoms. Comparing $\ell$ with its own smoothing avoids this.

\subsection{Gaussian smoothing}\label{sec:smoothing}

\begin{definition}[smoothing]\label{def:smoothing}
Let $L_\ell\colon\R\to\R$ be the piecewise-linear interpolant of the points $(\log k,\log\ell(k))$, $k\ge1$,
extended by the constant $\log\ell(1)$ on $(-\infty,0]$. Let $G(y):=e^{-y^2/2}/\sqrt{2\pi}$ and
\[
  \Lambda_\ell:=L_\ell*G,\qquad \rho_\ell(m):=\sup_{k\ge m}\bigl|\log\ell(k)-\Lambda_\ell(\log k)\bigr|\qquad(m\ge1).
\]
\end{definition}

The function $G$ extends to an entire function, and for $\xi=T+i\upsilon$ with $T,\upsilon,y\in\R$,
\begin{equation}\label{eq:gauss-abs}
  |G(\xi-y)|=G(T-y)\,e^{\upsilon^2/2}.
\end{equation}

\begin{lemma}[Gaussian smoothing]\label{lem:smoothing}
Let $\ell$ be slowly varying and write $L=L_\ell$, $\Lambda=\Lambda_\ell$.
\begin{enumerate}
\item[\textup{(o)}] For every $\eta>0$, $|L(y)|\le A_\eta+\eta(|y|+\log2)$ for all $y\in\R$, with $A_\eta$ as in
  Lemma~\ref{lem:potter}(ii). The function $\Lambda$ is entire and real on $\R$.
\item[\textup{(i)}] $\rho_\ell(m)\to0$ as $m\to\infty$.
\item[\textup{(ii)}] For every $\vartheta\ge0$, as $\Re\xi\to+\infty$, $\Lambda(\xi)-\Lambda(\Re\xi)\to0$ and
  $\Lambda^{(j)}(\xi)\to0$ for every $j\ge1$, uniformly in $|\Im\xi|\le\vartheta$; and $\Lambda$ is bounded on
  $\{\Re\xi\le0,\ |\Im\xi|\le\vartheta\}$.
\item[\textup{(iii)}] Consequently $\ell_1(k):=e^{\Lambda(\log k)}$ is analytic of angle $\vartheta$ for every
  $\vartheta>0$ \textup{(}Definition~\ref{def:analytic}, with any $c_0>0$\textup{)}, and $\psi:=\ell/\ell_1$ satisfies
  $\sup_{k\ge m}|\log\psi(k)|=\rho_\ell(m)$.
\end{enumerate}
\end{lemma}

The proof shows that $L$ has small oscillation on unit intervals far to the right, chains this into a bound for
$|L(y)-L(T)|$ that grows linearly in $|y-T|$, and integrates the bound against the derivatives of $G$, which decay
like a Gaussian even at complex arguments by \eqref{eq:gauss-abs}. We include the short proof, which uses only the
sequence form of slow variation.

\begin{proof}
\emph{(o).} For $y\ge0$ choose $k\ge1$ with $y\in[\log k,\log(k+1)]$. Then $L(y)$ is a convex combination of
$\log\ell(k)$ and $\log\ell(k+1)$, each of absolute value at most $A_\eta+\eta\log(k+1)\le A_\eta+\eta(y+\log2)$ by
Lemma~\ref{lem:potter}(ii). For $y<0$, $|L(y)|=|\log\ell(1)|\le A_\eta$. Since $L$ grows at most linearly, the integrand
$L(y)G(\xi-y)$ is dominated, for $\xi$ in a bounded set, by an integrable function of $y$, and it is entire in
$\xi$; so $\Lambda$ is entire and may be differentiated under the integral sign. It is real on $\R$ because $L$ and $G$
are.

\emph{Local oscillation of $L$.} For $T\ge T_0:=1+\log2$ put
$\osc_L(T):=\sup\{|L(y)-L(y')|:\ y,y'\in[T-1,T+1]\}$ and $\osc_L^*(T):=\sup_{T'\ge T}\osc_L(T')$. If
$y\in[T-1,T+1]\cap[\log k,\log(k+1)]$, then $k\ge e^{T-1}/2$ and $k+1\le e^{T+1}+1$. With $m:=\lceil e^T\rceil$ and
$\lambda:=2e^2$ we get $k,k+1\in[m/\lambda,\lambda m]$, so $\osc_L(T)\le2\osc_m(2e^2)$. Hence $\osc_L(T)\to0$, and
$\osc_L^*(T)$ is finite and tends to $0$.

\emph{Chaining.} Let $T\ge2T_0$. If $y\ge T$, join $T$ to $y$ by $\lfloor y-T\rfloor$ steps of length $1$ and one
shorter step; every step lies in an interval $[T'-1,T'+1]$ with $T'\ge T$, so $|L(y)-L(T)|\le\osc_L^*(T)(1+|y-T|)$.
If $T/2\le y<T$, the same argument started at $y\ge T_0$ gives $|L(y)-L(T)|\le\osc_L^*(T/2)(1+|y-T|)$. If $y<T/2$,
then $|y|+T\le5|y-T|$, and (o) with $\eta=1$ gives $|L(y)-L(T)|\le C_1(1+|y-T|)$ with $C_1:=\max(2A_1+2,5)$. In all
cases
\begin{equation}\label{eq:chaining}
  |L(y)-L(T)|\le\bigl(\osc_L^*(T/2)\,\mathbf 1_{\{y\ge T/2\}}+C_1\mathbf 1_{\{y<T/2\}}\bigr)(1+|y-T|).
\end{equation}

\emph{(ii).} The derivatives of $G$ are $G^{(j)}=\mathsf h_jG$ with polynomials $\mathsf h_j$ of degree $j$; let
$\bar{\mathsf h}_j$ have the absolute values of the coefficients of $\mathsf h_j$ as coefficients, so that
$|\mathsf h_j(t+i\upsilon)|\le\bar{\mathsf h}_j(|t|+\vartheta)$ for $|\upsilon|\le\vartheta$. The entire function
$\xi\mapsto\int_\R G(\xi-y)\diff y$ equals $1$ on $\R$, hence everywhere, and therefore
$\int_\R G^{(j)}(\xi-y)\diff y=\delta_{j0}$ for all $\xi$. For $\xi=T+i\upsilon$ with $|\upsilon|\le\vartheta$ and
$T\ge2T_0$, this, \eqref{eq:gauss-abs}, \eqref{eq:chaining} and the substitution $t=T-y$ give
\[
  \bigl|\Lambda^{(j)}(\xi)-\delta_{j0}L(T)\bigr|=\Bigl|\int_\R\bigl(L(y)-L(T)\bigr)G^{(j)}(\xi-y)\diff y\Bigr|
  \le e^{\vartheta^2/2}\Bigl[\osc_L^*(T/2)\int_\R\mathsf g_j+C_1\int_{|t|>T/2}\mathsf g_j\Bigr],
\]
where $\mathsf g_j(t):=\bar{\mathsf h}_j(|t|+\vartheta)(1+|t|)G(t)$ is integrable. The right side tends to $0$ as
$T\to\infty$, uniformly in $|\upsilon|\le\vartheta$. For $j=0$, applying this at $\xi$ and at $\Re\xi$ and subtracting
gives $\Lambda(\xi)-\Lambda(\Re\xi)\to0$; for $j\ge1$ it gives $\Lambda^{(j)}(\xi)\to0$. For $\Re\xi=T\le0$: $L(y)=L(0)$
for $y\le0$ and $|L(y)-L(0)|\le2A_1+2\log2+y$ for $y>0$, while $G(T-y)\le G(y)$ for $y>0$; hence
$|\Lambda(\xi)-L(0)|\le e^{\vartheta^2/2}\int_0^\infty(2A_1+2+y)G(y)\diff y$.

\emph{(i).} Since $L(\log k)=\log\ell(k)$, the case $j=0$, $\upsilon=0$ of the estimate in (ii) gives
$\Lambda(\log k)-\log\ell(k)\to0$.

\emph{(iii).} $\Lambda$ is entire, real on $\R$, bounded on $\{-c_0\le\Re\xi\le0,|\Im\xi|\le\vartheta\}$, and
$\sup\{|\Lambda'(\xi)|:\Re\xi\ge T,|\Im\xi|\le\vartheta\}\to0$, by (ii). Finally $\log\psi(k)=\log\ell(k)-\Lambda(\log k)$.
\end{proof}

\subsection{Stability under a factor tending to one}\label{sec:stability}
Fix $p>0$, $s>1$ and a slowly varying $\ell$, and let $\Lambda_\ell$, $\ell_1$, $\psi=\ell/\ell_1$ be as in
Lemma~\ref{lem:smoothing}. Put
\[
  \bar\rho:=\sup_k|\log\psi(k)|=\rho_\ell(1)<\infty,\qquad w^{(t)}_k:=k^{-s}\ell_1(k)\psi(k)^t\quad(t\in[0,1]),
\]
so that $w^{(0)}$ and $w^{(1)}$ are the weights of $\sigma_{p,s,\ell_1}$ and $\sigma_{p,s,\ell}$. Since
$\ell_1\psi^t=\ell\psi^{t-1}$, we have $e^{-\bar\rho}\ell(k)\le\ell_1(k)\psi(k)^t\le e^{\bar\rho}\ell(k)$. Write
$m^{(t)}_n$ for the spectral weights of $w^{(t)}$, call the indices $k<n/M_0$ the \emph{head}, and let
\[
  \mathfrak h_n:=\sup_{t\in[0,1]}\ \max_{n'\in\{n-1,n\}}\ m^{(t)}_{n'}\bigl([1,n/M_0)\bigr).
\]

\begin{lemma}[a priori size]\label{lem:apriori}
There are $C$ and $n_0$, depending on $p,s,\ell$, such that $\mathfrak h_n\le Cn^{s+5}e^{-2n}$ for all $n\ge2$, and
$a_n(w^{(t)})\le2M_0^pn^{-p}$ and $b_n(w^{(t)})\le2M_0^pn^{-p}$ for all $n\ge n_0$ and $t\in[0,1]$.
\end{lemma}

\begin{proof}
By Lemma~\ref{lem:potter}(ii) with $\eta=1$, $\|w^{(t)}\|_1\le e^{\bar\rho}\sigma_{p,s,\ell}(\R)$ and $w^{(t)}_k\ge
e^{-\bar\rho-A_1}k^{-s-1}$. For $n\ge2$ and $n'\in\{n-1,n\}$, the head $[1,n/M_0)$ is contained in $[1,2n'/M_0]$, so
Corollary~\ref{cor:head} with $A=s+1$ bounds $m^{(t)}_{n'}$ of the head by $Cn^{s+2}(n+1)^2e^{-2(n-1)}$, which gives the
first claim. For the second, $a_{n'}=\langle x,m_{n'}\rangle$ and $b_n\le\frac12\langle x,m_{n-1}+m_n\rangle$ by
\S\ref{sec:conv}. Off the head $x_k\le(M_0/n)^p$, and $x_k\le1$ always, so $a_n$ and $b_n$ are at most
$(M_0/n)^p+\mathfrak h_n$, which is at most $2M_0^pn^{-p}$ for large $n$.
\end{proof}

\begin{proposition}[stability]\label{prop:stability}
There are $C$ and $n_0$, depending on $p,s,\ell$, such that for $n\ge n_0$
\begin{align*}
  \bigl|\log b_n^2(\sigma_{p,s,\ell})-\log b_n^2(\sigma_{p,s,\ell_1})\bigr|
    &\le2\rho_\ell(\lceil n/M_0\rceil)+Cn^{s+5}e^{-2n},\\
  n^p\bigl|a_n(\sigma_{p,s,\ell})-a_n(\sigma_{p,s,\ell_1})\bigr|
    &\le6M_0^p\,\rho_\ell(\lceil n/M_0\rceil)+Cn^{p+s+5}e^{-2n}.
\end{align*}
\end{proposition}

\begin{proof}
Apply \S\ref{sec:fv} with $v_k=k^{-s}\ell_1(k)$ and $\varphi=\log\psi$, which is bounded by $\bar\rho$; then
$w(t)=w^{(t)}$. Write $\rho:=\rho_\ell(\lceil n/M_0\rceil)$. For $k\ge n/M_0$ we have $k\ge\lceil n/M_0\rceil$, so
$|\varphi(k)|\le\rho$; on the head $|\varphi(k)|\le\bar\rho$. By \eqref{eq:fv-b},
\[
  \bigl|\log b_n^2(w^{(1)})-\log b_n^2(w^{(0)})\bigr|\le\sup_{t\in[0,1]}\sum_k|\varphi(k)|\bigl(m^{(t)}_n(k)
  +m^{(t)}_{n-1}(k)\bigr)\le2\rho+2\bar\rho\,\mathfrak h_n .
\]
By \eqref{eq:fv-a} and $2|q_{n-1}q_n|\le m_{n-1}+m_n$,
$|\dot a_n|\le\sum_k|\varphi(k)|\bigl(|x_k-a_n|m_n(k)+b_n(m_{n-1}(k)+m_n(k))\bigr)$. Let $n\ge n_0$, with $n_0$ as in
Lemma~\ref{lem:apriori} and $2M_0^pn_0^{-p}\le1$. Off the head, $|x_k-a_n|\le2M_0^pn^{-p}$ and $b_n\le2M_0^pn^{-p}$, so
these indices contribute at most $6M_0^p\rho\,n^{-p}$; on the head use $|x_k-a_n|\le1$ and $b_n\le1$, a contribution of
at most $3\bar\rho\,\mathfrak h_n$. Integrate over $t\in[0,1]$ and apply Lemma~\ref{lem:apriori}.
\end{proof}

\subsection{Proof of Theorem~\ref{thm:sv}}\label{sec:proof-sv}
Let $\ell$ be slowly varying and $\ell_1:=e^{\Lambda_\ell(\log\cdot)}$. By Lemma~\ref{lem:smoothing}(iii), $\ell_1$ is
analytic of every angle, so Theorem~\ref{thm:imp-analytic} gives
$|n^pa_n(\sigma_{p,s,\ell_1})-2c|+|n^pb_n(\sigma_{p,s,\ell_1})-c|=O(1/n)$. Write $\rho:=\rho_\ell(\lceil n/M_0\rceil)$ and
$\Delta_n:=\log b_n^2(\sigma_{p,s,\ell})-\log b_n^2(\sigma_{p,s,\ell_1})$. By Proposition~\ref{prop:stability},
$|\Delta_n|\le2\rho+Cn^{s+5}e^{-2n}$ for $n\ge n_0$; in particular $\Delta_n$ is bounded, since $\rho\le\rho_\ell(1)$,
so $|e^{\Delta_n/2}-1|\le C'|\Delta_n|$ and
\[
  |n^pb_n(\sigma_{p,s,\ell})-c|\le|n^pb_n(\sigma_{p,s,\ell_1})-c|+n^pb_n(\sigma_{p,s,\ell_1})\,|e^{\Delta_n/2}-1|
  =O(1/n+\rho+n^{s+5}e^{-2n}).
\]
Similarly, the second bound of Proposition~\ref{prop:stability} gives
$|n^pa_n(\sigma_{p,s,\ell})-2c|\le|n^pa_n(\sigma_{p,s,\ell_1})-2c|+6M_0^p\rho+Cn^{p+s+5}e^{-2n}$. The exponentially
small terms are $O(1/n)$, and finitely many $n$ affect only the constant. This is \eqref{eq:sv-rate}, and the limits
follow because $\rho_\ell(m)\to0$ (Lemma~\ref{lem:smoothing}(i)). \qed

\begin{remark}[the order of the steps]\label{rem:order}
Theorem~\ref{thm:imp-analytic} is proved in \cite{PaperII} for every $s>1$, by an analysis in a window of exponents
followed by the Christoffel shift with rate \cite[Theorem~3.1]{PaperII}. The smoothing is applied here at the final
exponent $s$, after the shift. Smoothing first and shifting afterwards would also work, but the shift would then
have to carry the rate $1/n+\rho_\ell(\lceil n/M_0\rceil)$, which is in general not of the regular form required
there.
\end{remark}

\begin{remark}[what the rate measures]\label{rem:rate}
For $\ell$ that is itself analytic of some angle, \eqref{eq:sv-rate} is weaker than Theorem~\ref{thm:imp-analytic},
since $\rho_\ell$ need not be $O(1/n)$. Its point is the case of factors that are slowly varying but irregular at
every scale, such as $\ell(k)=1+\frac12(-1)^k/\log(k+1)$, for which $\rho_\ell(m)=O(1/\log m)$
(Section~\ref{sec:numerics}).
\end{remark}

% Section 4: rough weights (Theorem 1.3).
\section{Rough weights}\label{sec:rough}

This section proves Theorem~\ref{thm:rough}. Throughout, $s>1$ is fixed, $v_k:=k^{-s}$ is the pure power, a
superscript $v$ marks quantities of $\sigma_{p,s}$, and $w=ve^\varphi$ with $\varphi$ taking the values $0,\pm\varepsilon$.

The mechanism is the following. By Lemma~\ref{lem:lr}, $\log(b_N^2/b_{N-1}^2)$ and $a_N$ move, to first order, by
the pairing of $\varphi$ with $\Delta^2m_N^v$ and with $\Delta\cE_N^v$, up to a remainder of order $\varepsilon^2$ that holds
for every base measure. For the pure power these two vectors stay of size $\asymp1$ (after the natural scaling) on a
window of indices of size $N$; this follows from the arcsine limit of Lemma~\ref{lem:arcsine}, because a
degenerate vector would force the spectral weights to concentrate near three points
(Lemmas~\ref{lem:window} and~\ref{lem:akernel}). Choosing $\varphi=\pm\varepsilon$ with the sign of that vector on the
window produces a first-order response of size $\asymp\varepsilon$, which the remainder cannot cancel when
$\varepsilon$ is small. Repeating this on widely separated windows gives infinitely many such responses, and the
windows do not interact: the first-order term is evaluated at the pure power, where the other windows lie either in
the saturated head (Corollary~\ref{cor:head}) or beyond $Mn$, where the arcsine limit leaves little mass.

\subsection{Windows and the non-degeneracy of the second difference}\label{sec:window}
A \emph{window} is a pair $W=(\tau_1,\tau_2)$ with $0<\tau_1<\tau_2$; put $W_N:=\N\cap[\tau_1N,\tau_2N]$. It is
\emph{admissible} if $\tau_2>\usat(p)$. Let $Z:=\{-2,0,2\}$, and put
\begin{equation}\label{eq:window-constants}
\begin{gathered}
  I_W:=(\tau_2^{-p},\tau_1^{-p})\cap(0,4c),\qquad \Theta_W:=\{\theta\in(0,\pi):\ 2c(1+\cos\theta)\in I_W\},\\
  \iota_W:=\frac1\pi\int_{\Theta_W}\dist(2\cos\theta,Z)^2\diff\theta,\qquad \delta_W:=\frac{\iota_W}{32},\qquad
  J_W:=\Bigl\lceil\frac2{\delta_W}\Bigr\rceil,\qquad \gamma_W:=\frac{\delta_W}{J_W^2}.
\end{gathered}
\end{equation}
These depend only on $p,\tau_1,\tau_2$. The window is admissible if and only if $I_W\ne\emptyset$, since
$\tau_2^{-p}<4c$ is $\tau_2>\usat$. Then $\iota_W>0$: the map $\theta\mapsto2c(1+\cos\theta)$ takes $(0,\pi)$ onto
$(0,4c)$, so $\Theta_W$ is a non-empty open set, and the integrand vanishes only at $\theta\in\{0,\pi/2,\pi\}$. By
Lemma~\ref{lem:usat}, every window with $\tau_2\ge1$ is admissible.

Since $\arc_{4c}$ is the image of $\diff\theta/\pi$ under $y=2c(1+\cos\theta)$, and $(y-2c)/c=2\cos\theta$,
\begin{equation}\label{eq:iota-arc}
  \iota_W=\int_{I_W}\dist\bigl((y-2c)/c,Z\bigr)^2\diff\arc_{4c}(y).
\end{equation}

\begin{lemma}[window non-degeneracy]\label{lem:window}
Let $W$ be admissible. There is $N_0$ such that every $N\ge N_0$ has some $m\in[N-J_W+1,N]$ with
\[
  \bigl\|\Delta^2m^v_m|_{W_N}\bigr\|_1\ge\gamma_W .
\]
\end{lemma}

The proof is by contradiction. If the second differences are small on $W_N$ for $J_W$ consecutive indices, then so are
the first differences, so $|q_N|$, $|q_{N-1}|$ and $|q_{N-2}|$ nearly agree on $W_N$. The three-term recurrence then
forces the rescaled atoms in $W_N$, weighted by $m_{N-1}$, to lie near the three points $Z$, which contradicts the
arcsine limit, whose mass on $I_W$ is spread out.

\begin{proof}
Write $S:=W_N$, $\delta:=\delta_W$, $J:=J_W$, $\gamma:=\gamma_W$, $D_m:=m_m-m_{m-1}$, and drop the superscript $v$.
Suppose that $\|\Delta^2m_m|_S\|_1<\gamma$ for every $m\in[N-J+1,N]$.

\emph{Step 1: first differences are small on $S$.} For $m\in[N-J,N]$, $D_N-D_m=\sum_{j=m+1}^N\Delta^2m_j$, so
$\|(D_N-D_m)|_S\|_1\le(N-m)\gamma$. Also $m_N-m_{N-J}=\sum_{m=N-J+1}^ND_m$ has $\|\cdot|_S\|_1\le2$. Hence
$J\|D_N|_S\|_1\le2+\gamma J(J-1)/2$, that is, $\|D_N|_S\|_1\le2/J+(J-1)\gamma/2$, and since
$D_{N-1}=D_N-\Delta^2m_N$, also $\|D_{N-1}|_S\|_1\le2/J+(J+1)\gamma/2$. As $2/J\le\delta$ and $\gamma=\delta/J^2$, both
are at most $2\delta$.

\emph{Step 2: the moduli nearly agree on $S$.} Put $e:=|q_{N-2}|-|q_{N-1}|$ and $e':=|q_N|-|q_{N-1}|$. Since
$(\alpha-\beta)^2\le|\alpha^2-\beta^2|$ for $\alpha,\beta\ge0$, $\sum_{k\in S}e_k^2\le\|D_{N-1}|_S\|_1\le2\delta$ and
$\sum_{k\in S}e_k'^2\le\|D_N|_S\|_1\le2\delta$.

\emph{Step 3: the recurrence on $S$.} Fix $k\in S$, let $\epsilon_k:=\sgn(q_{N-1}(k)q_{N-2}(k))$ (with $\sgn0:=1$),
and $t_k:=x_k-a_{N-1}-b_{N-1}\epsilon_k$. We claim
\[
  |q_{N-1}(k)|\cdot\bigl||t_k|-b_N\bigr|\le b_{N-1}|e_k|+b_N|e'_k| .
\]
If $q_{N-1}(k)=0$ there is nothing to prove. Otherwise $q_{N-2}(k)=\epsilon_k\sgn(q_{N-1}(k))|q_{N-2}(k)|$, and
\eqref{eq:rec-vec} at index $N-1$, multiplied by $\sgn(q_{N-1}(k))$, gives
$b_N|q_N(k)|=\bigl|(x_k-a_{N-1})|q_{N-1}(k)|-b_{N-1}\epsilon_k|q_{N-2}(k)|\bigr|=\bigl|t_k|q_{N-1}(k)|-b_{N-1}\epsilon_ke_k\bigr|$.
Hence $\bigl|b_N|q_N(k)|-|t_k||q_{N-1}(k)|\bigr|\le b_{N-1}|e_k|$, and $b_N|q_N(k)|=b_N|q_{N-1}(k)|+b_Ne'_k$, which
proves the claim. Now $||t_k|-b_N|=\dist(t_k,\{\pm b_N\})\ge\dist(x_k-a_{N-1},\tilde Z)$ with
$\tilde Z:=\{\pm b_{N-1}\pm b_N\}$. Put $\bar b:=\max(b_N,b_{N-1})$, $y_k:=(x_k-a_{N-1})/\bar b$ and $Z_N:=\tilde
Z/\bar b$. Squaring the claim, using $(\alpha+\beta)^2\le2\alpha^2+2\beta^2$ and Step~2, and summing over $S$:
\begin{equation}\label{eq:concentration}
  \sum_{k\in S}m_{N-1}(k)\,\dist(y_k,Z_N)^2\le2\bigl(\|e|_S\|_2^2+\|e'|_S\|_2^2\bigr)\le8\delta=\iota_W/4 .
\end{equation}

\emph{Step 4: the arcsine limit contradicts \eqref{eq:concentration} for large $N$.} Put $n:=N-1$ and $Y_k:=n^px_k$.
Then $y_k=\phi_N(Y_k)$ with $\phi_N(Y):=(Y-n^pa_{N-1})/(n^p\bar b)$. By Theorem~\ref{thm:imp-pure},
$n^pa_{N-1}\to2c$ and $n^p\bar b\to c$, so $\phi_N\to\phi$, $\phi(Y):=(Y-2c)/c$, uniformly on compact sets, and
$Z_N\to Z$ in the Hausdorff distance. Since $\dist(\cdot,Z_N)$ is $1$-Lipschitz and differs from $\dist(\cdot,Z)$ by
at most the Hausdorff distance, $\dist(\phi_N,Z_N)^2\to\dist(\phi,Z)^2$ uniformly on compact sets. By
\eqref{eq:iota-arc} and monotone convergence there is a continuous $h$ with $0\le h\le1$, compact support in the open
interval $I_W$, and $\int h\,\dist(\phi,Z)^2\diff\arc_{4c}\ge\iota_W/2$. Now $k\in S$ if and only if $Y_k$ lies in
$[(n/(\tau_2N))^p,(n/(\tau_1N))^p]$, whose endpoints tend to $\tau_2^{-p}$ and $\tau_1^{-p}$; so for large $N$ the
support of $h$ lies inside this interval, and $\mathbf 1_S(k)\ge h(Y_k)$ for all $k$. Hence the left side of
\eqref{eq:concentration} is at least $\int G_N\diff\nu_{n,0}$ with $G_N:=h\,\dist(\phi_N,Z_N)^2$ and $\nu_{n,0}$ as in
Lemma~\ref{lem:arcsine}. The functions $G_N$ converge uniformly on $[0,\infty)$ to the bounded continuous function
$G:=h\,\dist(\phi,Z)^2$, so $\int G_N\diff\nu_{n,0}\to\int G\diff\arc_{4c}\ge\iota_W/2$ by Lemma~\ref{lem:arcsine}. For
$N$ large this contradicts \eqref{eq:concentration}.
\end{proof}

\subsection{The off-diagonal coefficients}\label{sec:rough-b}
Recall $F_N=\log(b_N^2/b_{N-1}^2)$. For an admissible window $W$ put
\[
  M_W:=\max\Bigl(1,\bigl(4c\sin^2(\pi\gamma_W/32)\bigr)^{-1/p}\Bigr),\qquad
  R_W:=\max\bigl(\tau_2/\tau_1,\ \tau_2M_0,\ M_W/\tau_1\bigr).
\]
Since $T_p$ of \eqref{eq:Tp} is non-increasing and $T_p((4c\sin^2(\pi\gamma_W/32))^{-1/p})=\gamma_W/16$, we have
$T_p(M_W)\le\gamma_W/16$.

\begin{theorem}\label{thm:rough-b}
Let $W$ be admissible, let $N_0$ be as in Lemma~\ref{lem:window}, and let $N_1<N_2<\cdots$ be integers with
$N_1\ge\max(N_0,2J_W+2)$, $\tau_2N_i<\tau_1N_{i+1}$, and either
\textup{(A)} $N_{i+1}/N_i\to\infty$, or \textup{(B)} $N_{i+1}>R_WN_i$ for every $i$. For each $i$ pick
$m_i\in[N_i-J_W+1,N_i]$ as in Lemma~\ref{lem:window}. Let $\varepsilon\in(0,\gamma_W/16]$ and
\[
  \varepsilon_k:=\varepsilon\,\sgn\bigl(\Delta^2m^v_{m_i}(k)\bigr)\quad(k\in W_{N_i},\ i\ge1),\qquad
  \varepsilon_k:=0\quad\text{otherwise},
\]
with $\sgn0:=1$. Then $w_k:=k^{-s}e^{\varepsilon_k}$ satisfies $\liminf_iF_{m_i}(w)\ge\varepsilon\gamma_W/2$ in case
\textup{(A)}, and $\liminf_iF_{m_i}(w)\ge\varepsilon\gamma_W/4$ in case \textup{(B)}.
\end{theorem}

\begin{proof}
The windows $W_{N_i}$ are disjoint, so $\varphi:=(\varepsilon_k)$ takes values in $\{0,\pm\varepsilon\}$. Fix $i$, write
$m:=m_i$, and let $H_i:=\bigcup_{j<i}W_{N_j}\subset[1,\tau_2N_{i-1}]$ and
$T_i:=\bigcup_{j>i}W_{N_j}\subset[\tau_1N_{i+1},\infty)$ be the other windows.

\emph{Linearisation.} Lemma~\ref{lem:lr}(ii), applied along the direct path $ve^{t\varphi}$, gives
$F_m(w)\ge F_m(v)+\langle\varphi,\Delta^2m^v_m\rangle-2\varepsilon^2$.

\emph{Own window and cross terms.} By the choice of the signs,
$\langle\varphi,\Delta^2m^v_m\rangle=\varepsilon\|\Delta^2m^v_m|_{W_{N_i}}\|_1+E_i$, where the first term is at least
$\varepsilon\gamma_W$ by Lemma~\ref{lem:window}, and the cross term $E_i:=\langle\varphi\mathbf 1_{H_i\cup T_i},
\Delta^2m_m^v\rangle$ satisfies
$|E_i|\le\varepsilon\bigl(m^v_m+2m^v_{m-1}+m^v_{m-2}\bigr)(H_i\cup T_i)$.

\emph{Head.} Let $m'\in\{m,m-1,m-2\}$; then $N_i-J_W-1\le m'\le N_i$, so $m'\ge N_i/2$. For $k\in H_i$,
$k/m'\le2\tau_2N_{i-1}/N_i$, which tends to $0$ in case (A) and is less than $2\tau_2/R_W\le2/M_0$ in case (B). So, for
$i$ large, $H_i\subset[1,2m'/M_0]$, and Corollary~\ref{cor:head} (with $c_1=1$, $A=s$, $\|v\|_1=\zeta(s)$) gives
$m^v_{m'}(H_i)\le\zeta(s)m'^{s+1}(m'+1)^2e^{-2m'}\to0$.

\emph{Tail.} In case (A) fix $M\ge1$. For $i$ large, $\tau_1N_{i+1}>MN_i\ge Mm'$, so $T_i\subset(Mm',\infty)$ and,
by Lemma~\ref{lem:arcsine}, $\limsup_im^v_{m'}(T_i)\le T_p(M)$; as $M$ is arbitrary, $m^v_{m'}(T_i)\to0$. In case (B),
$\tau_1N_{i+1}>\tau_1R_WN_i\ge M_Wm'$, so $\limsup_im^v_{m'}(T_i)\le T_p(M_W)\le\gamma_W/16$.

\emph{Conclusion.} Hence $\limsup_i|E_i|$ is $0$ in case (A) and at most $4\varepsilon\gamma_W/16$ in case (B). By
Theorem~\ref{thm:imp-pure}, $b_m(v)/b_{m-1}(v)\to1$, so $F_m(v)\to0$. With $\varepsilon\le\gamma_W/16$ we have
$2\varepsilon^2\le\varepsilon\gamma_W/8$, and therefore $\liminf_iF_{m_i}(w)$ is at least
$\varepsilon\gamma_W(1-\frac18)$ in case (A) and $\varepsilon\gamma_W(1-\frac18-\frac14)$ in case (B).
\end{proof}

\begin{proof}[Proof of Theorem~\ref{thm:rough}\textup{(i)}]
Integers $N_i$ as in Theorem~\ref{thm:rough-b} exist (choose them recursively), the weights constructed there are of
the stated form, and $\limsup_nF_n(w)\ge\liminf_iF_{m_i}(w)\ge\varepsilon\gamma_W/4$.
\end{proof}

\subsection{The diagonal coefficients}\label{sec:rough-a}
For the diagonal the relevant vector is $\Delta\cE_n$ of Lemma~\ref{lem:toda}, and its non-degeneracy has a direct proof,
with an explicit constant and a window of two indices. For $z\in\R$ put
\[
  \Upsilon(z):=\min_{\chi\in[0,\pi]}\bigl(|\sin2\chi-z\cos^2\chi|+|z\sin^2\chi-\sin2\chi|\bigr),
\]
and
\begin{equation}\label{eq:ia}
  \ia:=\frac1\pi\int_0^\pi\min\bigl(1,\Upsilon(2\cos\theta)\bigr)\diff\theta .
\end{equation}
The function $\Upsilon$ is $1$-Lipschitz, since the two terms change by at most $|z-z'|\cos^2\chi$ and
$|z-z'|\sin^2\chi$. It vanishes exactly on $Z=\{-2,0,2\}$: with $u=\cos\chi$, $u'=\sin\chi$ the two terms are
$|u(2u'-zu)|$ and $|u'(zu'-2u)|$; if $u=0$ the second is $|z|$, if $u'=0$ the first is $|z|$, and if $uu'\ne0$ both
vanish only if $2u'=zu$ and $zu'=2u$, whence $z^2=4$; conversely the minimum is $0$ for $z\in Z$, and it is attained.
Hence $0<\ia\le1$, because the integrand in \eqref{eq:ia} is continuous, non-negative, and vanishes only at
$\theta\in\{0,\pi/2,\pi\}$.

\begin{lemma}[non-degeneracy of the diagonal response]\label{lem:akernel}
For the pure power,
\[
  \liminf_{N\to\infty}\Bigl(\bigl\|N^p\Delta\cE^v_N\bigr\|_1+\bigl\|N^p\Delta\cE^v_{N+1}\bigr\|_1\Bigr)\ge2c\,\ia .
\]
Consequently, with $\gamma_{\mathrm a}=c\ia/2$, there is $N_0$ such that every $N\ge N_0$ has some $n\in\{N,N+1\}$ with
$\|n^p\Delta\cE^v_n\|_1\ge\gamma_{\mathrm a}$.
\end{lemma}

\begin{proof}
Fix $N$ and write $u:=q_N$, $u':=q_{N+1}$, $B:=b_{N+1}$ (all for $v$). By \eqref{eq:toda-id} at the indices $N$ and
$N+1$, at every atom
\[
  \Delta\cE_N=2Buu'-(x-a_N)u^2,\qquad \Delta\cE_{N+1}=(x-a_{N+1})u'^2-2Buu' .
\]
Put $y_k:=N^px_k$ and $z_k:=(y_k-2c)/c$. Then $N^p\Delta\cE_N=c(2uu'-zu^2)+r_N$ and
$N^p\Delta\cE_{N+1}=c(zu'^2-2uu')+r'_N$, where, by $2|uu'|\le u^2+u'^2$,
\[
  |r_N|+|r'_N|\le\delta_N(u^2+u'^2),\qquad \delta_N:=2|N^pB-c|+|N^pa_N-2c|+|N^pa_{N+1}-2c| ,
\]
and $\delta_N\to0$ by Theorem~\ref{thm:imp-pure}. The expression $|2uu'-zu^2|+|zu'^2-2uu'|$ is homogeneous of degree
$2$ in $(u,u')$, so it is at least $\Upsilon(z)(u^2+u'^2)$ at every atom. Summing over the atoms, with
$\sum_k(u_k^2+u_k'^2)=2$,
\[
  \|N^p\Delta\cE_N\|_1+\|N^p\Delta\cE_{N+1}\|_1\ge c\bigl\langle\min(1,\Upsilon(z)),m_N+m_{N+1}\bigr\rangle-2\delta_N .
\]
By Lemma~\ref{lem:arcsine} with $i=0$ and $i=1$, the images of $m_N$ and $m_{N+1}$ under $k\mapsto y_k$ converge weakly
to $\arc_{4c}$, and $y\mapsto\min(1,\Upsilon((y-2c)/c))$ is bounded and continuous; substituting $y=2c(1+\cos\theta)$,
the pairing tends to $2\ia$. This proves the first claim. For $N\ge N_0$ the sum is then at least $c\ia$, so one of the
two terms is at least $c\ia/2$; and $\|n^p\Delta\cE_n\|_1\ge\|N^p\Delta\cE_n\|_1$ for $n\ge N$.
\end{proof}

\begin{theorem}\label{thm:rough-a}
Let $\varepsilon\in(0,\ia/16]$. Let $N_0$ be as in Lemma~\ref{lem:akernel}, and choose integers $n_1<n_2<\cdots$ with
$\|n_i^p\Delta\cE^v_{n_i}\|_1\ge\gamma_{\mathrm a}$, $n_1\ge\max(N_0,2)$ and $n_{i+1}\ge(i+1)^2M_0n_i$. Put
$W_i:=\N\cap[n_i/M_0,\,in_i]$,
\[
  \varepsilon_k:=(-1)^i\varepsilon\,\sgn\bigl(\Delta\cE^v_{n_i}(k)\bigr)\quad(k\in W_i),\qquad\varepsilon_k:=0\quad
  \text{otherwise},
\]
and $w_k:=k^{-s}e^{\varepsilon_k}$. Then $(-1)^i\bigl(n_i^pa_{n_i}(w)-2c\bigr)\ge\varepsilon\gamma_{\mathrm a}/2-o(1)$ as
$i\to\infty$.
\end{theorem}

\begin{proof}
The indices $n_i$ exist, because every pair $\{N,N+1\}$ with $N\ge N_0$ contains one. The windows are disjoint: the
right end $in_i$ of $W_i$ is less than the left end $n_{i+1}/M_0\ge(i+1)^2n_i$ of $W_{i+1}$. Fix $i$ and put
$A_i:=\N\setminus W_i=[1,n_i/M_0)\cup(in_i,\infty)$. Every other window lies in $A_i$: for $j<i$,
$W_j\subset[1,(i-1)n_{i-1}]$ and $(i-1)n_{i-1}\le(i-1)n_i/(i^2M_0)<n_i/M_0$; for $j>i$, $W_j\subset[n_{i+1}/M_0,\infty)$
and $n_{i+1}/M_0>in_i$.

\emph{Off-window mass.} By Lemma~\ref{lem:toda} and Theorem~\ref{thm:imp-pure}, for $n$ large and every $A\subset\N$,
\[
  n^p\|\Delta\cE^v_n|_A\|_1\le\tfrac12n^p\max(b_n,b_{n+1})(v)\,\bigl(m^v_{n-1}+2m^v_n+m^v_{n+1}\bigr)(A)
  \le c\bigl(m^v_{n-1}+2m^v_n+m^v_{n+1}\bigr)(A).
\]
Let $\omega_i:=c(m^v_{n_i-1}+2m^v_{n_i}+m^v_{n_i+1})(A_i)$. We claim $\omega_i\to0$. For $n'\in\{n_i-1,n_i,n_i+1\}$,
the head $[1,n_i/M_0)$ lies in $[1,2n'/M_0]$, so Corollary~\ref{cor:head} makes its $m^v_{n'}$-mass exponentially
small. For the tail fix $M\ge1$: for $i\ge M$, $(in_i,\infty)\subset(Mn_i,\infty)$, and by Lemma~\ref{lem:arcsine}
with the shifts $-1,0,1$, $\limsup_im^v_{n'}((Mn_i,\infty))\le T_p(M)$. Hence $\limsup_i\omega_i\le4cT_p(M)$ for every
$M$, and $T_p(M)\to0$.

\emph{Linear response.} Lemma~\ref{lem:lr}(iii), applied along the direct path $ve^{t\varphi}$ with
$\varphi=(\varepsilon_k)$, gives
\[
  n_i^p\bigl|a_{n_i}(w)-a_{n_i}(v)-\langle\varphi,\Delta\cE^v_{n_i}\rangle\bigr|
  \le(e^\varepsilon-1-\varepsilon)\,n_i^p\bigl(b_{n_i}+b_{n_i+1}\bigr)(v)=(e^\varepsilon-1-\varepsilon)(2c+o(1)).
\]
By the choice of signs, and because $\varphi$ vanishes on $A_i$ except on the other windows,
\[
  (-1)^in_i^p\langle\varphi,\Delta\cE^v_{n_i}\rangle\ge\varepsilon\,n_i^p\|\Delta\cE^v_{n_i}|_{W_i}\|_1
  -\varepsilon\,n_i^p\|\Delta\cE^v_{n_i}|_{A_i}\|_1\ge\varepsilon(\gamma_{\mathrm a}-2\omega_i).
\]

\emph{Conclusion.} Since $\varepsilon\le\ia/16\le1/16$, $e^\varepsilon-1-\varepsilon\le\frac12\varepsilon^2e^\varepsilon\le
0.54\,\varepsilon^2$, and so $2c(e^\varepsilon-1-\varepsilon)\le1.08\,c\varepsilon\ia/16<\varepsilon\gamma_{\mathrm a}/2$.
With $n^pa_n(v)\to2c$ (Theorem~\ref{thm:imp-pure}) we obtain
$(-1)^i(n_i^pa_{n_i}(w)-2c)\ge\varepsilon\gamma_{\mathrm a}-\varepsilon\gamma_{\mathrm a}/2-o(1)$.
\end{proof}

\begin{proof}[Proof of Theorem~\ref{thm:rough}\textup{(ii)}]
The weights of Theorem~\ref{thm:rough-a} have the stated form, and the values at even and odd $i$ give the bounds for
the $\limsup$ and the $\liminf$.
\end{proof}

\subsection{The averaged law, and remarks}\label{sec:rough-rest}

\begin{proof}[Proof of Theorem~\ref{thm:rough}\textup{(iii)}]
If $e^{-\varepsilon}k^{-s}\le w_k\le e^{\varepsilon}k^{-s}$, then $\sum_kw_k\le e^\varepsilon\zeta(s)$ and $w_k\ge
e^{-\varepsilon}k^{-s}$, so Theorem~\ref{thm:imp-avg} gives \eqref{eq:paperI-logD} and \eqref{eq:paperI-cesaro}. In
case (i), suppose that $n^pb_n\to L\in[0,\infty]$. Then $\log(j^pb_j)\to\log L\in[-\infty,\infty]$, and the weights
$2(n-j)/n^2$ in \eqref{eq:paperI-cesaro} are non-negative with sum $(n-1)/n\to1$, so the left side of
\eqref{eq:paperI-cesaro} tends to $\log L$; hence $L=c$. But then $\log(b_n^2/b_{n-1}^2)=\log(\beta_n^2/\beta_{n-1}^2)
-2p\log\frac n{n-1}\to0$ with $\beta_n:=n^pb_n/c$, which contradicts~(i). In case (ii), the $\limsup$ and the
$\liminf$ of $n^pa_n$ differ by at least $\varepsilon\gamma_{\mathrm a}$.
\end{proof}

\begin{remark}[size of the constants]\label{rem:constants}
The constants of Theorem~\ref{thm:rough} are far from the size of the effect. For $p=2$ and $W=(0.6,1.6)$ a quadrature
gives $\iota_W\approx0.2345$, hence $J_W=273$, $\gamma_W\approx9.8\cdot10^{-8}$ and $R_W\approx1.1\cdot10^8$; for
$p=1$, $\iota_W\approx0.127$ and $\gamma_W\approx1.6\cdot10^{-8}$. Numerically, for the pure power with $p=s=2$, the
norm $\|\Delta^2m_m|_{W_N}\|_1$ is about $1.06$--$1.09$ for every $m\in[N-7,N]$ and $N\in\{100,200,300\}$, and
separation by a factor $4$ already suffices (Section~\ref{sec:numerics}). The loss comes from the contradiction in
Lemma~\ref{lem:window}, which costs the factor $J_W^2$. For the diagonal, a grid quadrature gives $\ia\approx0.314$,
and the observed first-order response $\varepsilon\|n^p\Delta\cE_n^v|_W\|_1$ is about $0.8\,\varepsilon c(p)$ for $p=s=2$,
about ten times the guaranteed $\varepsilon\gamma_{\mathrm a}/2$. None of these numerical values is used in the proofs,
which use only $\iota_W>0$ and $\ia>0$.
\end{remark}

\begin{remark}[the law to accuracy $\varepsilon^2$]\label{rem:accuracy}
Lemma~\ref{lem:lr}(i) shows what a perturbation of size $\varepsilon$ can do to the off-diagonal law. Let $\varphi$ be
bounded with $\frac12\osc\varphi\le\varepsilon$, $w=ve^\varphi$, and $\beta_n:=n^pb_n/c$. Then
$\log\beta_n^2(w)=\langle\varphi,m^v_n-m^v_{n-1}\rangle+r_n+\log\beta_n^2(v)$ with $|r_n|\le\varepsilon^2$ and
$\log\beta_n^2(v)\to0$, so
\[
  \Bigl|\limsup_n|\log\beta_n^2(w)|-\limsup_n|\langle\varphi,m^v_n-m^v_{n-1}\rangle|\Bigr|\le\varepsilon^2 .
\]
Replacing $\varphi$ by $\tau\varphi$ and letting $\tau\downarrow0$, the defect $\limsup_n|\log\beta_n^2(ve^{\tau\varphi})|$
is $o(\tau)$ if and only if $\langle\varphi,m^v_n-m^v_{n-1}\rangle\to0$. So a rough perturbation breaks the law at
first order exactly when it correlates with $m^v_n-m^v_{n-1}$ along infinitely many $n$, which is why the construction
uses sign patterns tied to the spectral weights. For fixed $\varphi$ this leaves a band of width $\varepsilon^2$ that we
cannot close; in particular we do not know whether $\langle\varphi,m_n^v-m_{n-1}^v\rangle\to0$ implies the law for
$ve^\varphi$.
\end{remark}

\begin{remark}[admissibility]\label{rem:admissible}
The condition $\tau_2>\usat$ is necessary for this mechanism. If $\tau_2\le\usat$ and $m-N$ stays bounded, the image
of $W_N$ under $k\mapsto m^px_k$ lies in $[(m/(\tau_2N))^p,\infty)$, whose left end tends to $\tau_2^{-p}\ge4c$; since
$\arc_{4c}$ has no mass beyond $4c$ and no atom, Lemma~\ref{lem:arcsine} gives $m^v_m(W_N)\to0$, so
$\|\Delta^2m^v_m|_{W_N}\|_1\to0$ and the first-order gain disappears. Windows of indices that are saturated in the
sense of Proposition~\ref{prop:imp-eq} carry no spectral weight.
\end{remark}

\begin{remark}\label{rem:rough-ab}
The construction of Theorem~\ref{thm:rough-a} controls $a_n$ only. We do not know whether $b_n$, or the ratio
$a_n/(2b_n)$, converges for the same weights.
\end{remark}

% Section 5: the Christoffel transformation at the Cesaro level (Theorem 1.4(i)).
\section{The Christoffel transformation at the Ces\`aro level}\label{sec:christoffel}

The direction $\varphi=-p\log k$ is special: $k^{-p}=x_k$, so $e^{-p\log k}v_k=x_kv_k$, and the end point of the path
$ve^{t\varphi}$ at $t=1$ is the measure $x\sigma$. For the pure powers this is the shift $s\mapsto s+p$. Here no
first-variation estimate is needed, because the Jacobi matrices of $\sigma$ and $x\sigma$ are related exactly by the
factorisation of Lemma~\ref{lem:imp-chain}. This section uses that factorisation beyond the pointwise transfer of
\cite{PaperII}: for the converse direction and a point mass at $0$ (Proposition~\ref{prop:converse}), and for the
product identity (Proposition~\ref{prop:product}), which with the chain-limit statement of
Lemma~\ref{lem:imp-chain}(ii) gives Theorem~\ref{thm:s}(i). Together with a comparison lemma for perturbations of
relative size $O(1/k)$ and two explicit measures, it determines the averaged $1/n$ coefficient on lattices of
exponents (Corollary~\ref{cor:lattice}). Finally it relates the pointwise $1/n$ coefficients at $s$ and $s+p$
(Proposition~\ref{prop:transfer}). In this section $\sigma$ is a positive Borel measure on $[0,\infty)$ with infinite
support and finite moments, $u_n,l_n$ are as in \eqref{eq:chain}, and $\tilde a_n,\tilde b_n$ are the coefficients of
$x\sigma$.

\subsection{The converse, and a point mass at zero}

\begin{proposition}\label{prop:converse}
Let $c'>0$. If $n^p\tilde a_n\to2c'$ and $n^p\tilde b_n\to c'$, then $n^pa_n\to2c'$ and $n^pb_n\to c'$. Consequently,
for every $m_0\ge0$, the measure $\sigma+m_0\delta_0$ satisfies $n^pa_n\to2c'$, $n^pb_n\to c'$ if and only if $\sigma$
does.
\end{proposition}

\begin{proof}
By Lemma~\ref{lem:imp-chain}(i), $l_{n+1}=\tilde a_n-u_n>0$ and $u_{n+1}=\tilde b_{n+1}^2/l_{n+1}$, so
$z_n:=n^pu_n/c'$ satisfies
\[
  z_{n+1}=\frac{B_n}{A_n-z_n},\qquad 0<z_n<A_n,\qquad A_n:=\frac{n^p\tilde a_n}{c'}\to2,\qquad
  B_n:=\frac{(n+1)^pn^p\tilde b_{n+1}^2}{c'^2}\to1 .
\]
The limiting map $z\mapsto1/(2-z)$ has the double fixed point $1$, and $1/(2-z)\ge z$ for $z<2$. Fix
$\delta\in(0,\frac15]$, put $\epsilon:=\delta^3/16$, and choose $n_5$ with $|A_n-2|,|B_n-1|\le\epsilon$ for $n\ge n_5$.
For $n\ge n_5$, $z_{n+1}<A_{n+1}$ gives $A_n-z_n>B_n/A_{n+1}\ge(1-\epsilon)/(2+\epsilon)>0.49$; hence $z_n<1.52$,
$2-z_n>0.48$, and
\[
  \Bigl|z_{n+1}-\frac1{2-z_n}\Bigr|\le\frac{|B_n-1|}{A_n-z_n}+\frac{|A_n-2|}{(A_n-z_n)(2-z_n)}\le7\epsilon .
\]
If $z_m\ge1+\delta$ for some $m\ge n_5$, then $z_{m+1}\ge z_m+(z_m-1)^2/(2-z_m)-7\epsilon\ge z_m+\delta^2/2$, and by
induction $z$ would exceed $1.52$; so $z_n<1+\delta$ for $n\ge n_5$. If $z_n\le1-\delta$, then
$z_{n+1}\ge z_n+\delta^2/4$, so $z$ enters $(1-\delta,1+\delta)$ after finitely many steps; and if $z_n>1-\delta$, then
$z_{n+1}\ge1/(1+\delta)-7\epsilon>1-\delta$. Hence $z_n\to1$. Finally $a_n=u_n+l_n=u_n+\tilde a_{n-1}-u_{n-1}$ and
$b_n^2=u_{n-1}l_n$ give $n^pa_n\to2c'$ and $n^{2p}b_n^2\to c'^2$.

For the second statement, $\sigma+m_0\delta_0$ has infinite support and finite moments, and
$x(\sigma+m_0\delta_0)=x\sigma$. By Theorem~\ref{thm:imp-shift} and the first part, the law for either measure is
equivalent to the law for $x\sigma$.
\end{proof}

The first statement also follows from the proposition of Dai, Piao and Qiao quoted in the introduction: with
$\mathsf g_n:=u_n/\tilde a_n\in(0,1)$ one has $\tilde b_n^2/(\tilde a_{n-1}\tilde a_n)=(1-\mathsf g_{n-1})\mathsf g_n$, so
these numbers form a chain sequence with parameter sequence $(\mathsf g_n)$, and the hypothesis makes them tend to
$\frac14$; by \cite[Proposition~2.1]{DPQ}, $\mathsf g_n\to\frac12$, which gives $n^pu_n\to c'$ and $n^pl_{n+1}\to c'$. For
discrete measures with bounded, strictly decreasing, non-negative atoms, the first statement of
Proposition~\ref{prop:converse} has been formalised in the Lean~4 proof assistant, with a general scale sequence in
place of $n^p$.

\subsection{The product identity and Theorem~\ref{thm:s}(i)}

\begin{proposition}\label{prop:product}
For every $n\ge1$,
\[
  \prod_{m=1}^n\Bigl(\frac{\tilde b_m}{b_m}\Bigr)^2=\frac{u_n}{u_0},\qquad u_0=a_0=\frac{\int x\diff\sigma}{\int\diff\sigma}.
\]
\end{proposition}

\begin{proof}
By Lemma~\ref{lem:imp-chain}(i), $b_m^2=u_{m-1}l_m$ and $\tilde b_m^2=u_ml_m$, so $(\tilde b_m/b_m)^2=u_m/u_{m-1}$;
the product telescopes. Also $u_0=-P_1(0)/P_0(0)=a_0$, since $P_1(x)=x-a_0$.
\end{proof}

\begin{proof}[Proof of Theorem~\ref{thm:s}\textup{(i)}]
Take $\sigma=\sigma_{p,s}$, so that $x\sigma=\sigma_{p,s+p}$ and $u_0=\zeta(s+p)/\zeta(s)$. Proposition~\ref{prop:product}
gives
\[
  2\sum_{m=1}^n\log\frac{b_m(s+p)}{b_m(s)}=\log\bigl(n^pu_n\bigr)-p\log n-\log\frac{\zeta(s+p)}{\zeta(s)} .
\]
By Theorem~\ref{thm:imp-pure}, $n^pa_n(s)\to2c$ and $n^pb_n(s)\to c$, so $n^pu_n\to c$ by
Lemma~\ref{lem:imp-chain}(ii).
\end{proof}

Only the law without rate is used. With the rate $O(1/n)$, the parameters converge only like
$n^pu_n-c=O(n^{-1/2})$ in general \cite[Section~3]{PaperII}, so the convergence in Theorem~\ref{thm:s}(i) has no
better rate a priori; numerically the error is of order $1/n$ (Section~\ref{sec:numerics}).

\subsection{Perturbations of relative size \texorpdfstring{$O(1/k)$}{O(1/k)}}

\begin{proposition}\label{prop:oneoverk}
Let $v,w$ be positive weights on the atoms $k^{-p}$ with $\sum_kv_k<\infty$, $v_k\ge c_1k^{-A}$ for some $c_1>0$,
$A\ge0$ and all $k$, and $|\log(w_k/v_k)|\le C/k$ for all $k\ge1$. Then $\log h_n(w)-\log h_n(v)=O(1/n)$, and
consequently
\[
  \sum_{j=1}^n\log\frac{b_j(w)}{b_j(v)}=-\frac12\log\frac{\sum_kw_k}{\sum_kv_k}+O(1/n).
\]
\end{proposition}

\begin{proof}
Let $\varphi:=\log(w/v)$ and $w(t):=ve^{t\varphi}$, $t\in[0,1]$; then $e^{-C}v\le w(t)\le e^Cv$. By \eqref{eq:fv-h},
$\frac{\diff}{\diff t}\log h_n=\langle\varphi,m^{(t)}_n\rangle$. For $k\ge n/M_0$, $|\varphi(k)|\le CM_0/n$. The head
$[1,n/M_0)$ lies in $[1,2n/M_0]$, so Corollary~\ref{cor:head}, with $\|w(t)\|_1\le e^C\|v\|_1$ and $w_k(t)\ge
e^{-C}c_1k^{-A}$, bounds its $m^{(t)}_n$-mass by $Cn^{A+3}e^{-2n}$, and there $|\varphi|\le C$. Hence
$|\frac{\diff}{\diff t}\log h_n|\le CM_0/n+O(n^{A+3}e^{-2n})$ uniformly in $t$; integrate over $[0,1]$. The second
statement follows from $h_n=h_0\prod_{j\le n}b_j^2$, where $h_0$ is the total mass.
\end{proof}

\subsection{Two anchors and the lattice values}
The averaged $1/n$ coefficient is fixed by an explicit measure at one exponent; Theorem~\ref{thm:s}(i) then moves it
along the lattice $s+p\N_0$. At $p=1$ the explicit measure is a Tricomi--Carlitz measure, which lives on a symmetric set
of atoms; the following standard lemma passes to the squared variable.

\begin{lemma}\label{lem:quadratic}
Let $\nu$ be a positive measure on $\R$, symmetric under $t\mapsto-t$, with infinite support and finite moments, and
let its monic orthogonal polynomials $R_N$ satisfy $tR_N=R_{N+1}+\mathsf e_NR_{N-1}$ \textup{(}$\mathsf e_0:=0$,
$\mathsf e_N>0$\textup{)}. Let $\sigma$ be the image of $\nu$ under $t\mapsto t^2$. Then $\sigma$ is a positive measure on
$[0,\infty)$ with infinite support and finite moments, its monic orthogonal polynomials $P_m$ satisfy
$R_{2m}(t)=P_m(t^2)$, and its coefficients are
\[
  a_m=\mathsf e_{2m}+\mathsf e_{2m+1}\quad(m\ge0),\qquad b_m^2=\mathsf e_{2m-1}\mathsf e_{2m}\quad(m\ge1).
\]
\end{lemma}

\begin{proof}
By symmetry $R_N(-t)=(-1)^NR_N(t)$, so $R_{2m}(t)=P_m(t^2)$ for a monic polynomial $P_m$ of degree $m$, and
$\int P_mx^j\diff\sigma=\int R_{2m}t^{2j}\diff\nu=0$ for $j<m$. So the $P_m$ are the monic orthogonal polynomials of
$\sigma$. Two steps of the recurrence give $t^2R_{2m}=R_{2m+2}+(\mathsf e_{2m}+\mathsf e_{2m+1})R_{2m}+\mathsf
e_{2m-1}\mathsf e_{2m}R_{2m-2}$, which is the recurrence of $P_m$ with the stated coefficients.
\end{proof}

\begin{corollary}[lattice values of the averaged $1/n$ coefficient]\label{cor:lattice}
\textup{(a)} For $p=1$ and every $s\in\frac32+\N_0$ there is $C(s)$ with
$\sum_{j=1}^n\log(2jb_j(s))=-\frac{s-1}2\log n+C(s)+o(1)$.

\textup{(b)} For $p=2$ and every $s\in2+2\N_0$ there is $C(s)$ with
$\sum_{j=1}^n\log(j^2b_j(s)/c(2))=-\frac{s-1}2\log n+C(s)+o(1)$.
\end{corollary}

\begin{proof}
(a) Let $\psi^{(1)}$ be the Tricomi--Carlitz measure with parameter $\alpha=1$: by \cite[(1.4)--(1.5)]{WLXZ} it is the
symmetric measure with mass $(k+1)^{k-1}e^{-k}/k!$ at each of the points $\pm(k+1)^{-1/2}$, $k\ge0$, and the
Tricomi--Carlitz polynomials $f_n=f^{(1)}_n$ are orthogonal with respect to it. By \cite[(1.9)]{WLXZ} their leading
coefficient is $\prod_{k<n}(k+1)/(k+1)=1$, so they are monic, and \cite[(1.6)]{WLXZ} reads
$(n+1)f_{n+1}(x)=(n+1)xf_n(x)-f_{n-1}(x)$, that is, $\mathsf e_N=1/(N+1)$ for $N\ge1$. The image of $\psi^{(1)}$ under
$t\mapsto t^2$ has mass $2(k+1)^{k-1}e^{-k}/k!=2e\,v_{k+1}$ at $1/(k+1)$, where $v_j:=j^{j-1}e^{-j}/j!$. By
Lemma~\ref{lem:quadratic}, the measure $\sigma_{\mathrm{TC}}:=\sum_{j\ge1}v_j\delta_{1/j}$ has the coefficients
$b^{\mathrm{TC}}_m=(2m(2m+1))^{-1/2}$.

Next compare $\sigma_{1,3/2}$ with $\sigma_{\mathrm{TC}}$. Put $\varpi(t):=\Gamma(t+1)e^t/(\sqrt{2\pi}\,t^{t+1/2})$, so that
$k^{-3/2}=\sqrt{2\pi}\,\varpi(k)v_k$. By Stirling's formula with the remainder bound for real argument \cite[(5.11.1) and
\S5.11(ii)]{DLMF}, $\log\Gamma(k)=(k-\frac12)\log k-k+\frac12\log(2\pi)+R(k)$ with $|R(k)|\le1/(12k)$; since
$\log\varpi(k)=\log\Gamma(k)+\log k+k-\frac12\log(2\pi)-(k+\frac12)\log k=R(k)$, we get $|\log\varpi(k)|\le1/(12k)$ for
all $k\ge1$. Apply Proposition~\ref{prop:oneoverk} with $w_k=k^{-3/2}$ and the weights $\sqrt{2\pi}\,v_k\ge
e^{-1/12}k^{-3/2}$, which have the same coefficients as $\sigma_{\mathrm{TC}}$:
\[
  \sum_{j=1}^n\log\frac{b_j(3/2)}{b^{\mathrm{TC}}_j}=-\frac12\log\frac{\zeta(3/2)}{\sqrt{2\pi}\,W_{\mathrm{TC}}}+O(1/n),
  \qquad W_{\mathrm{TC}}:=\sum_jv_j .
\]
Since $\log(2jb^{\mathrm{TC}}_j)=-\frac12\log(1+\frac1{2j})=-\frac1{4j}+O(j^{-2})$, this gives (a) for $s=\frac32$. By
Theorem~\ref{thm:s}(i) with $p=1$, $\sum_{j\le n}\log(b_j(s+1)/b_j(s))=-\frac12\log n+\mathrm{const}+o(1)$, so the
coefficient decreases by $\frac12$ at each step.

(b) By Proposition~\ref{prop:imp-lommel}, $j^2b_j(2)/c(2)=16j^2(4j+1)^{-1}((4j-1)(4j+3))^{-1/2}=1-\frac1{2j}+O(j^{-2})$, so
(b) holds for $s=2$; Theorem~\ref{thm:s}(i) with $p=2$ moves it by $-1$ at each step.
\end{proof}

\subsection{Transfer of the pointwise \texorpdfstring{$1/n$}{1/n} coefficient}
For $0<p\le1$ the pointwise $1/n$ term of $n^pb_n(s)$ is not known (Section~\ref{sec:open}). The product identity
shows how the coefficients at $s$ and $s+p$ must be related if they exist, and that their existence alone determines
the $1/n$ term of $a_n/(2b_n)$.

\begin{proposition}\label{prop:transfer}
Let $p>0$, $s>1$ and $\varepsilon>0$, and suppose that there are $d(s),d(s+p)\in\R$ with
\[
  \frac{n^pb_n(s)}{c}=1+\frac{d(s)}n+O(n^{-1-\varepsilon}),\qquad
  \frac{n^pb_n(s+p)}{c}=1+\frac{d(s+p)}n+O(n^{-1-\varepsilon}).
\]
Then $d(s+p)=d(s)-\frac p2$, and $n^pu_n(s)=c(1+O(n^{-\varepsilon'}))$ with $\varepsilon':=\min(\varepsilon,1)$. If moreover
$\varepsilon>\frac12$, then
\[
  \frac{a_n(s)}{2b_n(s)}=1-\frac p{2n}+O\bigl(n^{-2\varepsilon'}\bigr).
\]
\end{proposition}

\begin{proof}
By the hypothesis, $\log(b_m(s+p)/b_m(s))=(d(s+p)-d(s))/m+O(m^{-1-\varepsilon'})$, so the left side of
Theorem~\ref{thm:s}(i) is $(d(s+p)-d(s)+\frac p2)\log n+O(1)$. Its convergence forces $d(s+p)=d(s)-\frac p2$. Then,
by Proposition~\ref{prop:product}, $\log(u_m/u_{m-1})=2\log(b_m(s+p)/b_m(s))=-p/m+O(m^{-1-\varepsilon'})$, so
$\log(n^pu_n)$ converges with error $O(n^{-\varepsilon'})$, and its limit is $\log c$ by Lemma~\ref{lem:imp-chain}(ii).
For the last claim write $X:=b_n/u_{n-1}$ (at the exponent $s$). By Lemma~\ref{lem:imp-chain}(i),
$a_n=u_n+b_n^2/u_{n-1}$, so
\[
  \frac{a_n}{2b_n}=\frac12\Bigl(X+\frac{u_n/u_{n-1}}X\Bigr).
\]
Here $u_n/u_{n-1}=1-p/n+O(n^{-1-\varepsilon'})$, and $X=1+O(n^{-\varepsilon'})$, because
$u_{n-1}=c(n-1)^{-p}(1+O(n^{-\varepsilon'}))$ and $b_n=cn^{-p}(1+O(1/n))$. Hence
$a_n/(2b_n)=\frac12(X+X^{-1})-p/(2n)+O(n^{-1-\varepsilon'})=1+(X-1)^2/(2X)-p/(2n)+O(n^{-1-\varepsilon'})$, which is the
claim since $2\varepsilon'\le1+\varepsilon'$.
\end{proof}

For $p>1$ the hypothesis holds with $\varepsilon=1$ and $d(s)=-(s-1)/2$ by Theorem~\ref{thm:imp-second}, and the
conclusion agrees with that theorem. For $0<p\le1$ the proposition reduces the $1/n$ term of $a_n/(2b_n)$ to the
existence of the $1/n$ term of $n^pb_n$ at $s$ and at $s+p$, with a remainder $O(n^{-1-\varepsilon})$, $\varepsilon>\frac12$.

% Section 6: pure powers as s varies (Theorem 1.4(ii)-(iii), Corollary 1.5).
\section{Pure powers as \texorpdfstring{$s$}{s} varies}\label{sec:s}

This section proves Theorem~\ref{thm:s}(ii)--(iii) and Corollary~\ref{cor:p1}. Moving $s$ is the direction
$\varphi=-\log k$, which is unbounded, so we work with the exact derivative in $s$ (Proposition~\ref{prop:ds}) and with
a two-sided variational bound for $h_n(s)/h_n(s')$ (Lemma~\ref{lem:sandwich}). The bound involves a positive power of
$n^px$ under $m_n$, which the pointwise law controls, and a negative power, which is controlled near $x=0$ by the tail
bound of Lemma~\ref{lem:allM} (which rests on Proposition~\ref{prop:imp-tail}), at finitely many fixed exponents. Chaining over a grid of exponents gives
boundedness (Proposition~\ref{prop:UF}); the arcsine limit identifies the limit (Theorem~\ref{thm:UFlim}). This works
on compact subsets of the range $\Sp$ of the tail bound. The Christoffel identity of Theorem~\ref{thm:s}(i) then
removes the restriction to $\Sp$ for the averaged law (Lemma~\ref{lem:telescope}).

Throughout, $p>0$ is fixed, $m^{(s)}_n$ are the spectral weights of $\sigma_{p,s}$, $E_s$ denotes expectation with
respect to $m^{(s)}_n$ (with $n$ understood), $Y:=n^px_k=(n/k)^p$, and
\[
  \Phi_n(s):=\log h_n(s)+s\log n,\qquad \Theta_n(s):=\sum_k\log(k/n)\,m^{(s)}_n(k)\qquad(n\ge1).
\]
By \S\ref{sec:conv}, $E_s[Y]=n^pa_n(s)$.

\subsection{The derivative in \texorpdfstring{$s$}{s} and a variational bound}

\begin{proposition}\label{prop:ds}
For $n\ge0$ the function $s\mapsto h_n(s)$ is real-analytic on $(1,\infty)$, and
$\partial_s\log h_n(s)=-\sum_k(\log k)\,m^{(s)}_n(k)$. Hence $\Phi_n'(s)=-\Theta_n(s)$ for $n\ge1$.
\end{proposition}

\begin{proof}
The moments of $\sigma_{p,s}$ are $\zeta(s+pj)$, so $h_n=D_{n+1}/D_n$ is real-analytic, and so are the coefficients
of $P_n=P_n^{(s)}$. Differentiating $\sum_kk^{-s}P_n(x_k)^2$ term by term is justified because $|P_n(x_k)|\le
\max_{[0,1]}|P_n|$ and $\sum_kk^{-s}\log k$ converges locally uniformly; the term containing $\partial_sP_n$ vanishes,
since $\partial_sP_n$ has degree less than $n$. Hence $\partial_sh_n=-\sum_k(\log k)k^{-s}P_n(x_k)^2$.
\end{proof}

\begin{lemma}[variational bound]\label{lem:sandwich}
For $n\ge1$ and $s,s'>1$, with $E_{s'}$ taking values in $(0,\infty]$,
\[
  -\log E_{s'}\bigl[(k/n)^{s-s'}\bigr]\le\Phi_n(s)-\Phi_n(s')\le\log E_{s'}\bigl[(k/n)^{s'-s}\bigr].
\]
\end{lemma}

\begin{proof}
Write $P:=P_n^{(s')}$ and $\langle\cdot,\cdot\rangle_s$ for the inner product of $L^2(\sigma_{p,s})$. Upper bound: $h_n(s)$ is
the minimum of $\int Q^2\diff\sigma_{p,s}$ over monic $Q$ of degree $n$, so $h_n(s)\le\sum_kk^{-s}P(x_k)^2=
h_n(s')E_{s'}[k^{s'-s}]$; add $(s-s')\log n$ to the logarithms. Lower bound: assume $E_{s'}[k^{s-s'}]<\infty$ and put
$\phi(x_k):=k^{s-s'}P(x_k)$. Then $\|\phi\|_s^2=h_n(s')E_{s'}[k^{s-s'}]$ and $\langle Q,\phi\rangle_s=\langle
Q,P\rangle_{s'}$ for every polynomial $Q$, so $\phi$ is orthogonal in $L^2(\sigma_{p,s})$ to all polynomials of degree
less than $n$, and $\langle x^n,\phi\rangle_s=h_n(s')$. Writing $x^n=P_n^{(s)}+r$ with $\deg r<n$ and using the
Cauchy--Schwarz inequality, $h_n(s')\le h_n(s)^{1/2}\|\phi\|_s$, that is, $h_n(s)\ge h_n(s')/E_{s'}[k^{s-s'}]$.
\end{proof}

In terms of $Y$, $(k/n)^{s'-s}=Y^{(s-s')/p}$. For $s\ge s'$, the upper bound involves a positive power of $Y$, the
lower bound a negative one.

\subsection{Boundedness and the limit on compact subsets of \texorpdfstring{$\Sp$}{the window}}

\begin{proposition}\label{prop:UF}
For every compact interval $K=[\underline s,\bar s]\subset\Sp$, $\sup_{n\ge1}\sup_{s,s'\in K}|\Phi_n(s)-\Phi_n(s')|<\infty$.
\end{proposition}

\begin{proof}
Put $u_0:=\frac12\vartheta_*(p,\underline s)$, with $\vartheta_*$ as in Lemma~\ref{lem:allM}. Since $\vartheta_*$
increases with $s$, $u_0<\vartheta_*(p,s')$ for every $s'\in K$, and $u_0\le p/4$. Choose a grid $\underline s=t_0<t_1<\dots<t_J=\bar s$ with
$t_{i+1}-t_i\le u_0$. For each of the finitely many $t=t_i$, $A_t:=\sup_nn^pa_n(t)<\infty$ by
Theorem~\ref{thm:imp-pure}, and $B_t:=\sup_nE_t[(k/n)^{u_0}]<\infty$. For $n\ge n_0(t)$,
Lemma~\ref{lem:allM}, with exponent $\vartheta_0:=\vartheta_*(p,t)$, and Tonelli's theorem give
\[
  E_t[(k/n)^{u_0}]\le1+\sum_{k>n}m^{(t)}_n(k)\Bigl(1+\int_1^{k/n}u_0M^{u_0-1}\diff M\Bigr)
  \le2+\int_1^\infty u_0M^{u_0-1}CM^{-\vartheta_0}\diff M<\infty,
\]
and for the finitely many $n<n_0(t)$ each expectation is finite because $m^{(t)}_n(k)\le Ck^{-t}$ and $u_0<t-1$.

Let $s\in[t_i,t_{i+1}]$, $t:=t_i$ and $r:=(s-t)/p\in[0,\frac14]$. By Lemma~\ref{lem:sandwich} and Jensen's inequality
for the concave function $y\mapsto y^r$, $\Phi_n(s)-\Phi_n(t)\le\log E_t[Y^r]\le r\log E_t[Y]\le\frac14\log^+A_t$. By the
lower bound and $y^{pr}\le1+y^{u_0}$ for $y>0$ and $0\le pr\le u_0$, $\Phi_n(s)-\Phi_n(t)\ge-\log(1+B_t)$. So
$|\Phi_n(s)-\Phi_n(t_i)|$ is bounded on each cell uniformly in $n$, and chaining over the $J$ cells bounds
$|\Phi_n(s)-\Phi_n(\underline s)|$ on $K$.
\end{proof}

\begin{lemma}\label{lem:W}
Let $t\in\Sp$ and let $g$ be continuous on $(0,\infty)$ with $|g(y)|\le C(y^{-a}+y^{b})$, where $a\ge0$,
$pa<\vartheta_*(p,t)$ and $0\le b<2$. Then $E_t[g(Y)]\to\int g\diff\arc_{4c}$.
\end{lemma}

\begin{proof}
Choose $a'>a$ with $pa'<\vartheta_*(p,t)$, and put
\[
  V_-:=\sup_{n\ge n_0}E_t[Y^{-a'}],\qquad V_+:=\sup_nE_t[Y^2]=\sup_nn^{2p}\bigl(a_n^2+b_n^2+b_{n+1}^2\bigr)(t).
\]
As in the proof of Proposition~\ref{prop:UF}, $V_-<\infty$; and $V_+<\infty$ by Theorem~\ref{thm:imp-pure}, since
$E_t[x^2]=\langle x^2q_n,q_n\rangle$. The arcsine law satisfies the same two bounds, because its density is of order
$y^{-1/2}$ at $0$, $a'<\frac12$ (as $\vartheta_*<p/2$), and its support is compact. For $0<\epsilon<1<R$ the function
$g_{\epsilon,R}(y):=g(\min(\max(y,\epsilon),R))$ is bounded and continuous on $[0,\infty)$, and $\arc_{4c}$ has no atom at
$0$, so Lemma~\ref{lem:arcsine} gives $E_t[g_{\epsilon,R}(Y)]\to\int g_{\epsilon,R}\diff\arc_{4c}$. For each of these
measures,
\[
  E|g-g_{\epsilon,R}|\le E\bigl[(|g(Y)|+|g(\epsilon)|)\mathbf 1_{Y<\epsilon}\bigr]+E\bigl[(|g(Y)|+|g(R)|)\mathbf 1_{Y>R}\bigr]
  \le6CV_-\epsilon^{a'-a}+6CV_+R^{b-2},
\]
since $Y^{-a}$, $\epsilon^{-a}$ and $1$ are at most $Y^{-a'}\epsilon^{a'-a}$ on $\{Y<\epsilon\}$, and similarly on
$\{Y>R\}$. This tends to $0$ uniformly in $n$ as $\epsilon\to0$ and $R\to\infty$.
\end{proof}

Let $\phi(r):=\log\int y^r\diff\arc_{4c}(y)$ for $r>-\frac12$. It is finite, convex and $C^2$ on $[-\frac14,\frac14]$, with
$\phi(0)=0$, and there is $A_0$ with $|\phi(r)-r\phi'(0)|\le A_0r^2$ on $[-\frac14,\frac14]$. Moreover
\begin{equation}\label{eq:logmean}
  \phi'(0)=\int\log y\diff\arc_{4c}(y)=\log(2c)+\frac1\pi\int_0^\pi\log(1+\cos\theta)\diff\theta=\log c ,
\end{equation}
because $|1+e^{i\theta}|^2=2(1+\cos\theta)$ and $\frac1{2\pi}\int_0^{2\pi}\log|1+e^{i\theta}|\diff\theta=0$. The last
identity holds with $re^{i\theta}$, $r<1$, in place of $e^{i\theta}$ by the mean value property of the harmonic function
$\log|1+z|$, and passes to $r=1$ by dominated convergence, since $\log2\ge\log|1+re^{i\theta}|\ge\log|1+e^{i\theta}|-\frac12\log2$
for $\frac12\le r<1$ (use $|1+re^{i\theta}|^2\ge r|1+e^{i\theta}|^2$).

\begin{theorem}\label{thm:UFlim}
For every compact $K\subset\Sp$,
\[
  \sup_{s,s_0\in K}\Bigl|\Phi_n(s)-\Phi_n(s_0)-\frac{s-s_0}p\log c\Bigr|\longrightarrow0\qquad(n\to\infty).
\]
\end{theorem}

\begin{proof}
We may assume $K=[\underline s,\bar s]$. Let $u_0$ be as in the proof of Proposition~\ref{prop:UF}, let $0<u\le u_0$,
$r_u:=u/p\le\frac14$, and take the grid $t_i:=\underline s+iu$ in $K$ (the last cell may be shorter). For a cell
$[t_i,t_{i+1}]$ put $\phi_n(r):=\log E_{t_i}[Y^r]$; it is convex by H\"older's inequality, $\phi_n(0)=0$, and it is
finite for $|r|\le r_u$. For $s$ in the cell and $r:=(s-t_i)/p\in[0,r_u]$, Lemma~\ref{lem:sandwich} and convexity give
\[
  -\frac r{r_u}\phi_n(-r_u)\le-\phi_n(-r)\le\Phi_n(s)-\Phi_n(t_i)\le\phi_n(r)\le\frac r{r_u}\phi_n(r_u).
\]
By Lemma~\ref{lem:W} with $g(y)=y^{\pm r_u}$ (here $pr_u=u<\vartheta_*(p,t_i)$ and $r_u<2$),
$\phi_n(\pm r_u)\to\phi(\pm r_u)=\pm r_u\phi'(0)+e_\pm$ with $|e_\pm|\le A_0r_u^2$. Hence
$|\Phi_n(s)-\Phi_n(t_i)-r\phi'(0)|\le(r/r_u)(A_0r_u^2+\eta_{n,i})$ with $\eta_{n,i}\to0$, uniformly for $s$ in the cell.
Summing over the finitely many cells between $\underline s$ and $s$,
\[
  \limsup_{n\to\infty}\sup_{s\in K}\Bigl|\Phi_n(s)-\Phi_n(\underline s)-\frac{s-\underline s}p\phi'(0)\Bigr|
  \le A_0r_u\frac{\bar s-\underline s}p .
\]
As $u\in(0,u_0]$ is arbitrary, the left side is $0$, and $\phi'(0)=\log c$ by \eqref{eq:logmean}. Differences of two
points of $K$ follow.
\end{proof}

The grid has to be refined after $n\to\infty$; a fixed grid gives only the boundedness of Proposition~\ref{prop:UF}.
At $s-s_0=p$ the limit is that of Theorem~\ref{thm:s}(i), since $\Phi_n(s+p)-\Phi_n(s)=\log(n^pu_n(s))$ by
Proposition~\ref{prop:product}.

\subsection{Removing the restriction to \texorpdfstring{$\Sp$}{the window}}

\begin{lemma}\label{lem:telescope}
For all $s_0,s_1>1$, $\Phi_n(s_1)-\Phi_n(s_0)\to\frac{s_1-s_0}p\log c$.
\end{lemma}

\begin{proof}
For $s>1$ and $j\ge0$, Proposition~\ref{prop:product} and $h_n(s)=\zeta(s)\prod_{m\le n}b_m(s)^2$ give
$\Phi_n(s+(j+1)p)-\Phi_n(s+jp)=\log(n^pu_n(s+jp))$, which tends to $\log c$ by the proof of Theorem~\ref{thm:s}(i).
Write $s_i=t_i+j_ip$ with $t_i\in(1,1+p]$ and integers $j_i\ge0$ ($i=0,1$). Then
$\Phi_n(s_i)-\Phi_n(t_i)\to j_i\log c$. Since $(1,1+p]\subset\Sp$, Theorem~\ref{thm:UFlim} on the compact set
$\{t_0,t_1\}$ gives $\Phi_n(t_1)-\Phi_n(t_0)\to\frac{t_1-t_0}p\log c$. Adding,
$\Phi_n(s_1)-\Phi_n(s_0)\to(\frac{t_1-t_0}p+j_1-j_0)\log c=\frac{s_1-s_0}p\log c$.
\end{proof}

The lemma gives no uniformity outside $\Sp$: a uniform version would need the chain-limit statement of
Lemma~\ref{lem:imp-chain}(ii) uniformly in $s$, which we have not proved.

\subsection{Proofs of Theorem~\ref{thm:s}(ii)--(iii) and Corollary~\ref{cor:p1}}

\begin{proof}[Proof of Theorem~\ref{thm:s}\textup{(ii)}]
Since $h_n(s)=\zeta(s)\prod_{j\le n}b_j(s)^2$,
\[
  2\sum_{j=1}^n\log\frac{b_j(s_1)}{b_j(s_0)}=\Phi_n(s_1)-\Phi_n(s_0)-(s_1-s_0)\log n-\log\frac{\zeta(s_1)}{\zeta(s_0)} .
\]
Lemma~\ref{lem:telescope} gives the limit for all $s_0,s_1>1$, and Theorem~\ref{thm:UFlim} gives uniformity on compact
subsets of $\Sp$. The form with $h_n(s_1)/h_n(s_0)$ is the same statement.
\end{proof}

\begin{proof}[Proof of Corollary~\ref{cor:p1}]
Let $p=1$, so $c=\frac12$ and $\log(2jb_j(s))=\log(j^pb_j(s)/c)$. By Corollary~\ref{cor:lattice}(a) at $s_0=\frac32$,
$\sum_{j\le n}\log(2jb_j(\frac32))=-\frac14\log n+C(\frac32)+o(1)$, and by Theorem~\ref{thm:s}(ii),
$\sum_{j\le n}\log(b_j(s)/b_j(\frac32))=-\frac{s-3/2}2\log n+\mathrm{const}+o(1)$. Adding gives the first formula. The
moments of $\sigma_{1,s}$ are $\zeta(s+j)$, so $D_n(s)$ is its $n$-th Hankel determinant and
$D_{n+1}(s)/D_n(s)=h_n(s)=\zeta(s)\prod_{j\le n}b_j(s)^2$. Since
$2\sum_{j\le n}\log b_j=2\sum_{j\le n}\log(2jb_j)-2\log n!-2n\log2$, the second formula follows.
\end{proof}

\begin{proof}[Proof of Theorem~\ref{thm:s}\textup{(iii)}]
Let $K\subset\Sp$ be compact, fix $s_0\in K$ and put $\Delta_n(s):=\Phi_n(s)-\Phi_n(s_0)$. For $n\ge2$,
$\log b_n(s)^2=\log h_n(s)-\log h_{n-1}(s)$, so
\[
  \log\frac{b_n(s)^2}{b_n(s_0)^2}=\Delta_n(s)-\Delta_{n-1}(s)-(s-s_0)\log\frac n{n-1}.
\]
By Theorem~\ref{thm:UFlim}, $\Delta_n$ and $\Delta_{n-1}$ have the same uniform limit on $K$, so $b_n(s)/b_n(s_0)\to1$
uniformly on $K$; with $n^pb_n(s_0)\to c$ this is the first claim.

For the diagonal, Proposition~\ref{prop:product} gives $n^pu_n(s)=\exp(\Phi_n(s+p)-\Phi_n(s))$ whenever $s,s+p>1$, and
$a_n=u_n+b_n^2/u_{n-1}$ (Lemma~\ref{lem:imp-chain}(i)). If $p>1$, then $K\cup(K+p)$ lies in a compact subinterval of
$\Sp=(1,\infty)$, so $n^pu_n(s)\to c$ uniformly on $K$ by Theorem~\ref{thm:UFlim}, and each of $u_n$ and
$b_n^2/u_{n-1}$ is uniformly $cn^{-p}(1+o(1))$ on $K$. If $0<p\le1$ and $1+p\notin K$, write $K=K_1\cup K_2$ with the
compact sets $K_1:=K\cap(1,1+p)$ and $K_2:=K\cap(1+p,1+2p)$. On $K_1$ the previous argument applies, since
$K_1\cup(K_1+p)$ lies in a compact subinterval of $(1,1+2p)$. On $K_2$ use $\sigma_{p,s}=x\,\sigma_{p,s-p}$ and
Lemma~\ref{lem:imp-chain}(i) for $\sigma_{p,s-p}$:
\[
  a_n(s)=u_n(s-p)+l_{n+1}(s-p)=u_n(s-p)+\frac{b_{n+1}(s-p)^2}{u_n(s-p)} .
\]
Here $(K_2-p)\cup K_2$ lies in a compact subinterval of $(1,1+2p)$, so $n^pu_n(s-p)=\exp(\Phi_n(s)-\Phi_n(s-p))\to c$
uniformly on $K_2$ by Theorem~\ref{thm:UFlim}, and $(n+1)^pb_{n+1}(s-p)\to c$ uniformly on $K_2$ by the first claim
applied to $K_2-p$. Hence $n^pa_n(s)\to2c$ uniformly on $K_2$.
\end{proof}

At $s=1+p$ with $p\le1$ both arguments need $\Phi_n$ at an end point of the window ($s+p=1+2p$ or $s-p=1$). The
pointwise law holds there (Theorem~\ref{thm:imp-pure}), but we do not know whether the diagonal law is uniform near
$s=1+p$.

Theorem~\ref{thm:s}(iii) gives no rate. Numerically the deviations are of order $1/n$ (Section~\ref{sec:numerics}).

% Section 7: numerical illustrations.
\section{Numerical illustrations}\label{sec:numerics}

The computations below are floating-point illustrations; no proof uses them. The scripts and their stored outputs
are in the repository named in the data availability statement. Coefficients were computed either by the Lanczos
process with full reorthogonalisation in double precision on the atoms $k\le K$, or, where a tail correction is
mentioned, from the atoms $k\le K$ together with a Gauss--Jacobi rule for the continuous tail
$\int_{K+1/2}^\infty t^{-s}g(t^{-p})\diff t$, by an algorithm that adds one atom at a time and restores tridiagonal form by
plane rotations, in quadruple precision; the latter reproduces the Lommel coefficients of
Proposition~\ref{prop:imp-lommel} to a relative error below $10^{-14}$ for $n\le400$.

\emph{Rough weights.} Table~\ref{tab:rough} shows the construction of Theorem~\ref{thm:rough-b} for $p=s=2$, the window
$W=(0.6,1.6)$, three scales $N_i=15,60,240$ (separation by the factor $4$, far below $R_W$) and two values of
$\varepsilon$, both far above the proved range $\varepsilon\le\gamma_W/16\approx6\cdot10^{-9}$. At each scale the jump
$F_m(w)-F_m(v)$ at the selected index $m$ agrees with the first-order prediction $\langle\varphi,\Delta^2m_m^v\rangle$
to within $2\%$, the contribution of the other windows is at most $1.4\cdot10^{-3}$, and every jump exceeds the
largest $|F_n(w)-F_n(v)|$ over the indices $n\in[10,246]$ with $|n-m_i|>2$ for all $i$ by a factor larger than $2$. The measure was truncated at
$K=9600$ atoms; with two scales $(30,240)$ and $\varepsilon=0.2$, doubling $K$ changes the jumps by less than
$10^{-3}$.

\begin{table}[ht]
\centering
\small
\begin{tabular}{rrrrrr}
\toprule
$\varepsilon$ & $N_i$ & $m_i$ & $F_{m_i}(w)-F_{m_i}(v)$ & $\langle\varphi,\Delta^2m^v_{m_i}\rangle$ & other windows\\
\midrule
0.05 & 15  & 15  & $0.0689$ & $0.0691$ & $-3.5\cdot10^{-4}$\\
0.05 & 60  & 58  & $0.0624$ & $0.0626$ & $2.0\cdot10^{-4}$\\
0.05 & 240 & 238 & $0.0601$ & $0.0601$ & $<10^{-16}$\\
0.2  & 15  & 15  & $0.2739$ & $0.2765$ & $-1.4\cdot10^{-3}$\\
0.2  & 60  & 58  & $0.2457$ & $0.2503$ & $8.2\cdot10^{-4}$\\
0.2  & 240 & 238 & $0.2389$ & $0.2403$ & $<10^{-16}$\\
\bottomrule
\end{tabular}
\caption{The construction of Theorem~\ref{thm:rough-b} for $p=s=2$, $W=(0.6,1.6)$, scales $15,60,240$. The largest
$|F_n(w)-F_n(v)|$ over $n\in[10,246]$ with $|n-m_i|>2$ is $0.026$ for $\varepsilon=0.05$ and $0.103$ for
$\varepsilon=0.2$. Double precision, $K=9600$ atoms.}
\label{tab:rough}
\end{table}

\emph{The Christoffel--Ces\`aro identity and the $s$-law.} Table~\ref{tab:shift} lists the error
$S_n:=\sum_{m\le n}\log(b_m(s+p)/b_m(s))+\frac p2\log n-\frac12\log(c\zeta(s)/\zeta(s+p))$ in Theorem~\ref{thm:s}(i),
computed with the tail correction ($K=3.2\cdot10^5$ atoms, and $K=6.4\cdot10^5$ for $p=\frac12$). It decays like $1/n$, much faster than the
a priori rate $n^{-1/2}$. For the uniform statements, on the grid $s\in\{1.5,1.75,\dots,4\}$ and with $K=4\cdot10^4$ plus the
tail correction, the quantity $\sup_s|\Phi_n(s)-\Phi_n(\frac32)-(s-\frac32)\log c/p|$ of Theorem~\ref{thm:UFlim} equals
$2.25\cdot10^{-2}$, $1.12\cdot10^{-2}$, $5.62\cdot10^{-3}$, $2.81\cdot10^{-3}$ for $p=1$ and $5.26\cdot10^{-2}$,
$2.64\cdot10^{-2}$, $1.32\cdot10^{-2}$, $6.63\cdot10^{-3}$ for $p=2$, at $n=50,100,200,400$; the deviations
$\sup_s|n^pb_n(s)/c-1|$ and $\sup_s|n^pa_n(s)/(2c)-1|$ are $3.7\cdot10^{-3}$ and $5.0\cdot10^{-3}$ ($p=1$), and
$3.7\cdot10^{-3}$ and $6.2\cdot10^{-3}$ ($p=2$), at $n=400$, and halve with each doubling of $n$. For $p=1$ the grid extends
beyond $\Sp=(1,3)$; there the computation illustrates what we expect but have not proved. Finally, for the
Tricomi--Carlitz comparison in the proof of Corollary~\ref{cor:lattice}(a), the error of
$\sum_{j\le n}\log(b_j(\frac32)/b^{\mathrm{TC}}_j)$, computed with the value $W_{\mathrm{TC}}=1$ that the computation
indicates, is $4.1\cdot10^{-4}$ at $n=100$ and $2.6\cdot10^{-5}$ at $n=1600$.

\begin{table}[ht]
\centering
\small
\begin{tabular}{rrrrrr}
\toprule
$p$ & $s$ & $n=100$ & $n=400$ & $n=800$ & $n=1599$\\
\midrule
0.5 & 1.5 & $-6.4\cdot10^{-5}$ & $-6.7\cdot10^{-6}$ & $-2.2\cdot10^{-6}$ & $-7.6\cdot10^{-7}$\\
1   & 2   & $-2.5\cdot10^{-3}$ & $-6.3\cdot10^{-4}$ & $-3.1\cdot10^{-4}$ & $-1.6\cdot10^{-4}$\\
1   & 3   & $2.3\cdot10^{-3}$  & $6.1\cdot10^{-4}$  & $3.1\cdot10^{-4}$  & $1.6\cdot10^{-4}$\\
1.5 & 1.5 & $-8.3\cdot10^{-3}$ & $-2.1\cdot10^{-3}$ & $-1.0\cdot10^{-3}$ & $-5.2\cdot10^{-4}$\\
2   & 2   & $-9.9\cdot10^{-3}$ & $-2.5\cdot10^{-3}$ & $-1.2\cdot10^{-3}$ & $-6.3\cdot10^{-4}$\\
\bottomrule
\end{tabular}
\caption{The error $S_n$ in the Christoffel--Ces\`aro identity, Theorem~\ref{thm:s}(i). Quadruple precision, atoms
plus a quadrature of the continuous tail.}
\label{tab:shift}
\end{table}

\emph{Slowly varying factors.} For $p=s=2$ and the unbounded factor $\ell(k)=\exp\sqrt{\log(k+e)}$, the relative
deviation $n^pb_n(\sigma_{p,s,\ell})/(n^pb_n(\sigma_{p,s}))-1$ equals $1.09\cdot10^{-3}$, $5.23\cdot10^{-4}$,
$2.49\cdot10^{-4}$, $1.18\cdot10^{-4}$, $5.6\cdot10^{-5}$ at $n=100,200,400,800,1600$ (extended precision, $K=1.2\cdot10^5$
and $2.4\cdot10^5$ atoms without tail correction, which agree to $2\cdot10^{-7}$); $n$ times it decreases slowly, like
$\frac12\frac{\diff}{\diff\xi}\sqrt{\log(e^\xi+e)}$ at $\xi=\log n$. For the factor $\ell(k)=1+\frac12(-1)^k/\log(k+1)$ of
Remark~\ref{rem:rate}, which is slowly varying but rough at every scale, the deviation $|b_n(\sigma_{2,2,\ell})/b_n(\sigma_{2,2})-1|$
is $3.9\cdot10^{-3}$, $1.5\cdot10^{-3}$ and $8.2\cdot10^{-4}$ at $n=250,1000,3600$ (extended precision, a single truncation
at $K=1.2\cdot10^5$ atoms, so these values are only indicative); the decay is slow, as the rate $\rho_\ell(\lceil
n/M_0\rceil)=O(1/\log n)$ of Theorem~\ref{thm:sv} allows.

% Section 8: open problems and conjectures.
\section{Open problems}\label{sec:open}

\emph{Between slow variation and boundedness.} Theorems~\ref{thm:sv} and~\ref{thm:rough} leave a range of regularity
conditions open. A simple test case is $w_k=k^{-s}e^{\varepsilon\sin\log k}$: the increments of $\varepsilon\sin\log k$
tend to $0$, but $e^{\varepsilon\sin\log k}$ is not slowly varying, so Theorem~\ref{thm:sv} does not apply, and the
first-order response of Remark~\ref{rem:accuracy} has not been analysed for it. For bounded rough weights one may also
ask whether the pointwise law holds along a subsequence of density one; we do not know. Averaged forms of the diagonal
law, such as Ces\`aro means of $a_n$ over tails, are the natural counterpart of \eqref{eq:paperI-cesaro} for the
diagonal; we do not treat them here.

\emph{Uniformity in $s$.} For $0<p\le1$, Theorem~\ref{thm:s}(ii)--(iii) are uniform only on compact subsets of
$(1,1+2p)$, and the diagonal part of (iii) only away from $s=1+p$. Both restrictions come from the range of the tail
bound, Proposition~\ref{prop:imp-tail}, and would disappear with a uniform version of the chain-limit statement of
Lemma~\ref{lem:imp-chain}(ii). A related question is the size of the logarithmic moment $\Theta_n(s)=-\Phi_n'(s)$.
Theorem~\ref{thm:UFlim} shows that its integrals over intervals of $s$ converge, and for each fixed $s\in\Sp$ the bound of
Lemma~\ref{lem:allM} (for $k>n$) and the law $n^pa_n\to2c$ (for $k<n$, via $\log(n/k)\le(n/k)^p/p$) give
$\Theta_n(s)=O(1)$; whether $\sup_{s\in K}|\Theta_n(s)|=O(1)$ on compact sets is open. It
would follow from the next conjecture, because a lower bound on $\Phi_n''=-\Theta_n'$ bounds $\Theta_n$ on $K$ by its values
at one point.

\begin{conjecture}[convexity in $s$]\label{conj:convex}
For every $p>0$ and every compact $K\subset(1,\infty)$ there is $n_0$ such that $s\mapsto\log h_n(\sigma_{p,s})$ is
convex on $K$ for $n\ge n_0$.
\end{conjecture}

The evidence is numerical: second differences in $s$ of $\log h_n$, from exact determinants for $n\le8$ and from the
computed coefficients at $n=1600$ on the grid $s\in\{\frac32,2,3,4\}$, $p\in\{\frac12,1,\frac32,2,3,4\}$, are positive,
of size about $1/(p^2n)$. So the conjecture concerns the $1/n$ level and is not a soft statement. It also uses more than
the exponential-family structure $\sigma_{p,s}=x^{(s-s')/p}\sigma_{p,s'}$, as the following remark shows.

\begin{remark}\label{rem:cvx-fin}
Let $f_{ij}:=\int\tilde P_i\tilde P_j\log k\diff\sigma$, where $\tilde P_j$ are the orthonormal polynomials of a measure
$\sigma=\sum_kk^{-s}\delta_{k^{-p}}$ with finitely or infinitely many atoms and $\log k$ is evaluated at the atom $k^{-p}$.
Differentiating $\int\tilde P_n\tilde P_i\diff\sigma=\delta_{ni}$ in $s$ as in Lemma~\ref{lem:fv} gives
$\partial_s^2\log h_n=\Var_{m_n}(\log k)-2\sum_{i<n}f_{in}^2$. For the truncation $\sigma^{(N)}:=\sum_{k\le N}k^{-s}\delta_{k^{-p}}$
and $n=N-1$, the space $L^2(\sigma^{(N)})$ is spanned by $\tilde P_0,\dots,\tilde P_{N-1}$, so
$\Var_{m_{N-1}}(\log k)=\sum_{i<N-1}f_{i,N-1}^2$ and $\partial_s^2\log h_{N-1}=-\sum_{i<N-1}f_{i,N-1}^2$. This is negative:
otherwise $(\log k)\tilde P_{N-1}$ would be a multiple of $\tilde P_{N-1}$ on the atoms, so $\tilde P_{N-1}$ would vanish at
all atoms but one, which is impossible since it is orthogonal to the constants. So convexity fails at the top level of
every finite truncation.
\end{remark}

\emph{The pointwise $1/n$ term for $p\le1$.} For $p>1$ the $1/n$ term is Theorem~\ref{thm:imp-second}. For $0<p\le1$,
Corollary~\ref{cor:p1} gives its logarithmic average at $p=1$, Theorem~\ref{thm:s}(ii) gives its dependence on $s$ in the
same averaged sense for every $p$, and Proposition~\ref{prop:transfer} reduces the diagonal ratio to the off-diagonal
term.

\begin{conjecture}[{\cite[Conjecture~11.1]{PaperII}}]\label{conj:pleq1}
For $0<p\le1$ and $s>1$, $n^pb_n(\sigma_{p,s})=c(p)\bigl(1-\frac{s-1}{2n}+o(1/n)\bigr)$ and
$a_n(\sigma_{p,s})/(2b_n(\sigma_{p,s}))=1-\frac p{2n}+o(1/n)$.
\end{conjecture}

Numerically, at $p=\frac12$ the remainder in the first formula appears to be of order $n^{-1-p}=n^{-3/2}$, which is exactly
the borderline $\varepsilon=\frac12$ of Proposition~\ref{prop:transfer}; so even a proof of the first formula with this
remainder would not, through Proposition~\ref{prop:transfer} alone, give the second.

\emph{Slowly varying factors at order $1/n$.} The numbers in Section~\ref{sec:numerics} suggest that a slowly varying
factor acts at order $1/n$ like a change of $s$ by its local logarithmic derivative. This is the pointwise analogue of
Theorem~\ref{thm:s}(ii) along the path $k^{-s}\ell(k)^t$.

\begin{conjecture}[{\cite[Conjecture~11.2]{PaperII}}]\label{conj:tilt}
Let $\ell$ be analytic of some angle, $\ell(k)=\exp\Lambda(\log k)$, and $s>1$. If $n\Lambda'(\log n)\to\infty$, then
\[
  \frac{n^pb_n(\sigma_{p,s,\ell})}{c(p)}=1-\frac{s-1-\Lambda'(\log n)}{2n}+o\Bigl(\frac{\Lambda'(\log n)}n\Bigr).
\]
\end{conjecture}

\emph{Other atoms.} All results here use the exact atoms $k^{-p}$: Lemma~\ref{lem:arcsine} needs a pointwise law for a
reference measure, and Section~\ref{sec:christoffel} needs $x\sigma_{p,s}=\sigma_{p,s+p}$. For atoms $x_k=f(k)$ regularly
varying of index $-p$ the averaged law is \cite[Theorem~1.1]{PaperI}, and the pointwise law is open.


\section*{Acknowledgements}
The research and the exposition were developed with the assistance of an AI system (Claude, Anthropic). The main
theorem of the first paper of this series is formally verified in the Lean proof assistant. The author takes
responsibility for the content.

\section*{Data and code availability}
The scripts that produced the numbers of Section~\ref{sec:numerics}, their stored outputs, and the Lean sources
mentioned in Section~\ref{sec:christoffel} (declarations \texttt{christoffel\_shift\_law\_converse} and
\texttt{chain\_limit}) are in the repository \url{https://github.com/speagleg/jacobi-asymptotics} (tag
\texttt{v1.0-paper3} [placeholder: tag to be created at submission]).

\bibliographystyle{amsplain}
\bibliography{refs}

\end{document}
